% before: gal-local-situation
Both halves of this chapter compare representations of $\operatorname{GL}_n$ of a local field with
representations of a Weil group. The Weil group is taken inside the absolute Galois group of a number
field, as the subgroup of a decomposition group acting on the residue field by integral powers of
Frobenius. This lets the same group receive restrictions of global Galois representations.

% before: gal-local-artin
A local Langlands correspondence matches characters of $E_w^\times$ with characters of the Weil
group, and the matching is made through a reciprocity map. The normalization chosen sends geometric
Frobenius elements to uniformizers. Compatibility with the global reciprocity map ties the local map
to the global Galois group in which the Weil group sits.

% before: gal-local-artin-unique
The laws below are stated with respect to a local Artin map, and two maps satisfying them are
compared later. That comparison needs the local Artin map to be unique. Existence and uniqueness are
local class field theory.

% before: gal-weil-deligne
The objects on the Galois side of the correspondence are Weil--Deligne representations with
semisimple Frobenius action. Isomorphism classes are what the correspondence produces, so sums,
twists, tensor products and duals are defined on classes. The special representation
$\operatorname{Sp}(l)$ carries the monodromy of a segment of length $l$.

% before: gal-wd-classes
Local-global compatibility will be asserted only for certain kinds of Weil--Deligne
representations. Abelian representations are those of characters of the Weil group, tame ones are
trivial on wild inertia, and the multiplicity condition singles out tame representations whose
repeated inertial eigenvalues come from characters of $k_w^\times$. These are the classes on which
the correspondence below is determined.

% before: gal-wd-galois
An $\ell$-adic representation of $G_E$ restricts to the Weil group at a place prime to $\ell$, and
Grothendieck's monodromy theorem turns the restriction into a Weil--Deligne representation. The
definition is arranged so that the class is always defined. For continuous representations it is the
usual one.

% before: gal-p-adic-hodge
At places over $\ell$ the comparison with automorphic data goes through $p$-adic Hodge theory. The
Galois representations of this chapter are crystalline at unramified places over $\ell$, and at some
other places they are potentially semistable with a Weil--Deligne representation given by Fontaine's
functor. Their Hodge--Tate weights are read off from the Hodge filtration.

% before: gal-irreducible-classes
The automorphic side of the correspondence is the set of irreducible smooth representations of all
the groups $\operatorname{GL}_n(E_w)$. The laws are stated through the Bernstein--Zelevinsky
supercuspidal support and the Langlands classification by segments. Depth zero is the condition
under which a supercuspidal is described by a cuspidal representation of a finite general linear
group.

% before: gal-good
The correspondence used in the proof is determined only on part of
$\operatorname{Irr}(E_w)$. A supercuspidal is good if it is a character, or if $E_w$ is
$\mathbf Q_g$ and the supercuspidal is a twist of one of depth zero. Good representations are those
built from good supercuspidals.

% before: gal-green-datum
On depth-zero supercuspidals over $\mathbf Q_g$ the correspondence is fixed by Green's
parametrization of cuspidal representations of $\operatorname{GL}_n(\mathbf F_g)$. Comparing a Green
parameter with a Galois representation needs the residue field of degree $n$ over $\mathbf F_g$ and
the reduction of a reciprocity map. A Green datum supplies both.

% before: gal-llc-laws
The proof needs a map from $\operatorname{Irr}(E_w)$ to Weil--Deligne representations with
enough properties to read off the local behaviour of Galois representations, without the full local
Langlands correspondence.
The laws below hold for every $\pi$ when they concern the dimension, characters, the Langlands
classification and absolute values, and only for good $\pi$ when they concern twists, central
characters, duals, injectivity and tameness. The last law describes depth-zero supercuspidals over
$\mathbf Q_g$ through their Green parameters.

% before: gal-llc-exists
Every later statement quantifies over maps satisfying the laws. Their existence is needed so that
these statements apply to at least one map at every place. The map is Alp\"oge's construction of
Section~\ref{section-gal-local-construction}: local class field theory on characters, explicit tame
parameters on twists of depth-zero supercuspidals over some $\mathbf Q_g$, and placeholders elsewhere.

% before: gal-depth-zero-parameter
The laws determine the value of a map at a good supercuspidal of depth zero through a small list of
properties. This lemma says that these properties single out one Weil--Deligne representation, and
that this representation determines the supercuspidal. It is the local input for the uniqueness
statement that follows.

% before: gal-llc-unique
Two maps satisfying the laws need not agree on all of $\operatorname{Irr}(E_w)$. They do agree on
good representations, because the laws pin down their values on characters, on twists of depth-zero
supercuspidals, and then through the Langlands classification. This is the sense in which the laws
define a local Langlands correspondence on good representations.

% before: gal-llc-characterization
To relate the laws to the published correspondence we list the properties that characterize it.
These include the functional equation of Jacquet, Piatetski-Shapiro and Shalika with Deligne's local
constants, bijectivity, and the description of spherical and depth-zero representations.

% before: gal-llc-of-rec
The characterization implies every law except the description of depth-zero supercuspidals through
Green parameters. The next statement takes that description as a hypothesis. It says that a map with
the published characterization satisfies the laws whenever the hypothesis holds.

% before: gal-llc-scope
The laws are weaker than the published characterization in two ways. They contain no $L$-factors or
$\varepsilon$-factors, and they say little about representations that are not good. A comparison with
the published correspondence is available only under the hypothesis of the last proposition.

% before: gal-depth-zero-facts
The construction of a map satisfying the laws rests on the theory of depth-zero representations of
$\operatorname{GL}_n(\mathbf Q_g)$. A depth-zero supercuspidal is determined by the cuspidal
representation of $\operatorname{GL}_n(\mathbf F_g)$ it contains and by its central character. These
facts, and the effect of twists on depth, are used in each step below.

% before: gal-tame-inertia-character
Green's parametrization labels cuspidal representations of $\operatorname{GL}_n(\mathbf F_g)$ by
characters of $\mathbf F_{g^n}^\times$, and these must be transported to characters of inertia. The
residue of the local Artin map of an unramified extension of degree $n$ does this. Its properties,
and its independence of the auxiliary number field, are what the type law needs.

% before: gal-induced-parameter
An irreducible tame Weil--Deligne representation of dimension $n$ is induced from a character of an
unramified extension of degree $n$. The lemma builds it directly from the $n$ conjugate characters of
inertia and a prescribed determinant of Frobenius. It is used twice: for parameters of depth-zero
supercuspidals and for the representations assigned to supercuspidals that are not good.

% before: gal-depth-zero-parameter-exists
Lemma~\ref{gal-depth-zero-parameter} shows that a parameter is unique and determines the
supercuspidal. The other half is existence. Green's character of the cuspidal type gives the inertial
characters, and the central character fixes the determinant.

% before: gal-parameter-operations
The map on good supercuspidals must commute with twists and contragredients. Since parameters are
unique, it is enough to check that twisting and dualizing a parameter give a parameter of the twisted
or dual representation. Twists are by tame characters, so that depth zero is preserved.

% before: gal-good-map
A good supercuspidal is a character or a twist of a depth-zero supercuspidal, and in general in more
than one way. The value on it is defined through any such description. Uniqueness of parameters
shows that the descriptions give the same value.

% before: gal-placeholder
On supercuspidals that are not good the laws ask little, but injectivity at good representations
forbids their values from coinciding with values at good ones. The lemma provides, in every dimension
$n\geq2$, an irreducible representation that is wildly ramified in a way no twist of a tame one is.

% before: gal-exponent
The placeholder of Lemma~\ref{gal-placeholder} has determinant of absolute value $1$, while the laws
require the absolute value of the central character. A real exponent attached to each
representation corrects this by an unramified twist. It shifts by $s$ under $\nu^s$, so that the law
on unramified twists holds as well.

% before: gal-cusp-laws
The construction first defines the map on supercuspidals and then extends it through the Langlands
classification. The laws split accordingly. This definition lists the part that concerns
supercuspidals.

% before: gal-cusp-map
The map on supercuspidals is assembled from the map on good supercuspidals and the twisted
placeholders. Each law on supercuspidals is then checked from the lemmas above, with injectivity
using that the placeholder is not a twist of a tame representation.

% before: gal-langlands-extension
Every irreducible representation is the Langlands quotient of a unique multiset of segments. Law~(3)
therefore leaves no freedom once the map is known on supercuspidals. The extension is defined by that
formula.

% before: gal-extension-laws
The laws on supercuspidals now have to be propagated to all irreducible representations. The
supercuspidal support and the segments of a representation determine its central character, twists
and contragredient, and goodness and depth zero are read off the support.

% before: gal-polarized
The Galois representations are attached to cuspidal representations of $\operatorname{GL}_a$ over a
CM field that are regular algebraic and polarized. Regular algebraicity fixes the Hodge--Tate
weights, and polarization relates $\pi$ to its conjugate and its contragredient. These are the
hypotheses of the theorem of Barnet-Lamb, Gee, Geraghty and Taylor.

% before: gal-compatible
Local-global compatibility at a place prime to $\ell$ compares the Weil--Deligne representation of
the Galois representation with the image of the local component under a map satisfying the laws.
Here the comparison is asked only when that image is abelian and monodromy-free, or tame,
monodromy-free and multiplicity-free over some $\mathbf Q_g$ with $g$ large. In the second case only
the monodromy and the inertial characteristic polynomials are compared.

% before: gal-attached
The properties asked of a Galois representation are collected for a class of maps $\mathcal R$ at
each place. The weights are given by the infinitesimal character, and the comparison at places over
$\ell$ uses the same restricted cases as away from $\ell$. No polarization or purity is part of the
list.

% before: gal-galois-representations
Later chapters use the existence of attached Galois representations for every polarized regular
algebraic cuspidal representation. The statement is phrased for a class of maps $\mathcal R$ so that
it can be applied with the maps satisfying the laws.

% before: gal-galois-representations-any-rc
This is the form of the existence of Galois representations used in Chapters~\ref{chapter-tate-shafarevich}
and~\ref{chapter-assembly}. Its local-global compatibility holds for every map satisfying the laws,
so it applies to the map of Theorem~\ref{gal-llc-exists}. The construction is Alp\"oge's, from the
ingredients of the next section, and is carried out in Sections~\ref{section-gal-seed-construction}
and~\ref{section-gal-construction}.

% before: gal-galois-scope
The conditions in the definition of attached triples are stated on the values of a map satisfying
the laws. By the laws, they translate into conditions on the local component itself; for instance, an
abelian monodromy-free value comes from a local component whose supercuspidal support consists of
characters. Outside these cases the definition gives no information.

% before: gal-cyclic-base-change
The Galois representations are built after base change to solvable extensions of the CM field. Cyclic
base change of prime degree transfers cuspidal representations of $\operatorname{GL}_a$ upward, and
it keeps them cuspidal unless they are invariant under twisting by a character of the extension.

% before: gal-definite-unitary
The construction passes through automorphic representations of the definite unitary group of the
standard Hermitian form. Their local components at split places are representations of general
linear groups, and their weights at infinity are highest weights of compact unitary groups. These
are the data that base change and theta lifting relate.

% before: gal-base-change-definite
To return from definite unitary groups to general linear groups the construction uses Labesse's
base change. It sends an automorphic representation of $U_n$ to an isobaric sum of Speh
representations of conjugate self-dual cuspidal representations, compatibly with local components at
split places, unramified places and infinity.

% before: gal-good-base
Theta lifting between unitary groups requires a choice of additive character, a CM type and lattices
in the Hermitian and skew-Hermitian spaces. A good base records these choices together with
conditions at the inert places. The two kinds of spaces are those used in the theta correspondence
below; each also carries a Cartan involution at every infinite place, which fixes the Gaussian vectors
of the Weil representation there, and a basis of its lattice at the split places of the base.

% before: gal-theta-lift-large-rank
The seed representations exist only for unitary groups of odd rank with lacunary weights. Theta lifts
to unitary groups of large rank carry an automorphic representation of $U_n$ to one of $U_d$ for
every large $d$, with explicit weights at infinity and with the local theta correspondence at split
places. This moves automorphic representations into the range where seeds exist. Nonvanishing of
theta lifts to unitary groups of large rank goes back to Li \cite{Li90}; see also \cite{HKS96}.

% before: gal-seed
Seed representations are the first Galois representations of the construction. They may have large
dimension, and their Frobenius eigenvalues at almost all suitable places are only known to lie among
the normalized Satake parameters. A lacunary set of Hodge--Tate weights is what allows a piece of the
right dimension to be cut out later.

% before: gal-seed-representations
The seeds come from the cohomology of PEL unitary Shimura varieties, with Frobenius eigenvalues
controlled by a congruence relation and Hodge--Tate weights computed from a holomorphic Lefschetz
formula. The hypotheses fix a type at one place over $\ell$ and lacunary weights at another. The
conclusions give a seed in general, seeds adapted to a given place prime to $\ell$, and crystalline
seeds at a given place over $\ell$. The statement is Alp\"oge's; its geometric ingredients go back to
Kottwitz \cite{Kot92}, Harris and Taylor \cite{HT01}, and Faltings and Chai \cite{FC90}.

% before: gal-hecke-tower
The seeds of Theorem~\ref{gal-seed-representations} are pieces of the \'etale cohomology of smooth
projective varieties cut out by Hecke correspondences. A Hecke tower records the variety and the
correspondences. Its Galois pieces and its Hodge cohomology carry the information that the seeds
need.

% before: gal-tower-reduction
Frobenius eigenvalues of a piece at a place of good reduction are controlled by a congruence relation
between the Frobenius correspondence and Hecke correspondences. The definition records the smooth
model on which such a relation is stated and the Hecke polynomial whose roots it constrains.

% before: gal-seed-situation
The statements of this section share the hypotheses of Theorem~\ref{gal-seed-representations}: a CM
field with enough splitting, an automorphic representation of a definite unitary group with a
supercuspidal type at one place over $\ell$, and a second place over $\ell$. The lacunary weight is
added where it is used.

% before: gal-pack
A pack collects the geometric data that produce a seed for $U$: a Hecke tower, a Hecke eigensystem
on it, and an algebraic Hecke character that adjusts the Frobenius eigenvalues. The congruence
condition links the eigensystem to the Satake parameters of $U$ at almost all $d$-split places. Hodge
data at a place over $\ell$ say which Hodge degrees the eigensystem occupies.

% before: gal-etale-foundations
The cohomology of a Hecke tower enters through two theorems of \'etale cohomology. Finiteness makes the
Hecke eigenspaces finite-dimensional Galois representations, and smooth and proper base change
compares the cohomology of the generic and special fibres of a model, which is where the Frobenius
elements at a place of good reduction can be computed.

% before: gal-padic-comparison
At places over $\ell$ the seeds must be crystalline with prescribed Hodge--Tate weights and Frobenius.
These come from Faltings's comparison theorems for smooth projective schemes over $p$-adic integers,
used in the three forms stated here: crystallinity, the Hodge--Tate decomposition, and the identification
of the crystalline Frobenius.

% before: gal-weighted-trace
A finite correspondence with multiplicities acts on cohomology through a trace along its finite leg,
weighted by the multiplicities. The properties of this weighted trace listed here are the ones used to
compare such operators: additivity in the weight, compatibility with pushing the weight forward, and
the fact that a weight is determined by its values over a dense set.

% before: gal-weighted-correspondences
The Eichler--Shimura relation is an identity between finite correspondences that holds on points of the
special fibre. To use it on cohomology one needs that a finite correspondence with multiplicities acts on
cohomology through its $0$-cycles at geometric points, and that these may be tested on a dense set of
points.

% before: gal-eichler-shimura-cohomology
With the specialization isomorphism and the previous lemma, the Eichler--Shimura relation at the generic
points of the special fibre becomes a polynomial identity between the geometric Frobenius and the Hecke
operators on the cohomology of the generic fibre. Inertia acts trivially there because the model is
smooth and proper.

% before: gal-crystalline-test-algebras
At a place over $\ell$ the Eichler--Shimura relation is an identity between points of the special
fibre in characteristic $\ell$, while Proposition~\ref{gal-tower-cohomology}(5) needs an identity for
the crystalline Frobenius on the cohomology of the generic fibre. Alpöge connects the two through the
cohomology of functors of points on perfect algebras over the tilt of $\mathcal O_C$, with coefficients
in their crystalline rings. On points with values in $A/\ell^\flat$ the $q$-power Frobenius is a natural
transformation, and it can be compared there with $\phi^f$.

% before: gal-passage-ordered-rings
A test algebra $A$ is absolutely integrally closed and $\ell^\flat$-ordered in the sense defined next,
and the arguments below need the same two properties for its untilt $A^\sharp$, with an $\ell$-th root
of $\ell$ in place of $\ell^\flat$. Alpöge transfers them to the quotient of $A^\sharp$ by the square of
that root and then lifts them to $A^\sharp$ by successive approximation. The lifting is done in the
quotients $S/I_P$, where every idempotent of $S$ becomes $0$ or $1$, and Newton steps are its
approximation step.

% before: gal-passage-ordered-basics
Items (1) and (2) carry the two defining properties of a test algebra $A$ to quotients of $A$, among
them the quotient of $A^\sharp$ by the square of an $\ell$-th root of $\ell$. Items (4) and (5) give the
hypotheses of the splitting of \'etale covers below, both for $A^\sharp$ and for $A^\sharp[1/\ell]$.

% before: gal-passage-integral-closure
In a $\varpi$-ordered ring, divisibility by a power of $\varpi$ can be tested on the two pieces cut out
by an ordering idempotent, on each of which one of two elements divides the other. This gives the
integral closedness of the untilt and the extraction of roots from divisibilities of $d$-th powers that
the successive approximation below needs.

% before: gal-passage-components
An idempotent $e\notin P$ becomes $1$ in $S/I_P$, so an identity in $S/I_P$ holds in $S$ after
multiplication by such an idempotent. Item (2) assembles such local statements over all maximal ideals
into a global one. Items (3) and (4) compare the $\varpi$-ordered property of $S$ with the chained
property of $\varpi$ in the rings $S/I_{\mathfrak m}$.

% before: gal-passage-chained-mod-square
That $S/\varpi^2S$ is ordered controls divisibility in $S$ only up to multiples of $\varpi^2$. When
$\varpi$ lies in the Jacobson radical this suffices to compare two elements that generate an ideal
containing a power of $\varpi$, because a divisor of a power of $\varpi$ is a power of $\varpi$ times a
divisor of $\varpi$. Item (2) supplies the hypothesis of item (1) in the rings $S/I_P$.

% before: gal-passage-newton
Alpöge finds roots of monic polynomials in $S$ by raising, one step at a time, the power of $\varpi$
that divides the value. For one step he substitutes $cX$ for $X$, with $c$ read off from the
divisibilities among the coefficients as for the first slope of a Newton polygon, so that the constant
term divides the other coefficients. A root modulo $\varpi^2$ of the quotient by the constant term then
gives the step.

% before: gal-passage-newton-limit
Newton steps are found in the rings $S/I_{\mathfrak m}$, where $\varpi$ is chained, and item (1)
assembles them into Newton steps in $S$. In item (2) a Newton step bounds only the $d$-th power of the
correction, so the convergence of the successive approximations uses the extraction of roots of
Lemma~\ref{gal-passage-integral-closure}(1).

% before: gal-passage-lift-square
For a $\varpi$-adically complete ring $S$, the ordered property of $S/\varpi^2S$ lifts to $S$ through
the rings $S/I_P$, and absolute integral closedness lifts through Newton steps. This applies to
$S=A^\sharp$ for a test algebra $A$, with $\varpi$ an $\ell$-th root of $\ell$, because
$A^\sharp/\varpi^2$ is a quotient of $A$.

% before: gal-passage-etale-sections
Item (2) of Lemma~\ref{gal-crystalline-test-facts} asks for sections of \'etale surjections onto
$\operatorname{Spec}A^\sharp$ and $\operatorname{Spec}A^\sharp[1/\ell]$. Over an absolutely integrally
closed ring a standard \'etale algebra has a point near a maximal ideal as soon as a root of its
defining polynomial avoids the inverted element there. The componentwise local property spreads that
point to an open and closed neighbourhood, and Lemma~\ref{gal-passage-components}(2) glues these
neighbourhoods.

% before: gal-passage-untilt
The untilt $A^\sharp=W(A)/\xi$ is $\ell$-adically complete and $\ell$-torsion free, and modulo $\ell$
it is $A/\ell^\flat$. Since $A$ is perfect, $\ell^\flat$ has an $\ell$-th root in $A$, which gives an
$\ell$-th root $\varpi$ of $\ell$ in $A^\sharp$. Then $A^\sharp/\varpi^2$ is a quotient of $A$, so the
two defining properties of $A$ pass to it.

% before: gal-crystalline-test-facts
The arguments below split finite \'etale covers over the untilts of test algebras and refine open
covers by decompositions. Both are possible because a test algebra behaves like a product of valuation
rings with algebraically closed fraction fields.

% before: gal-shifting-coefficient-systems
Classes with coefficients in $A_{\mathrm{cris}}/\ell^\nu$ are studied by changing the coefficients:
embedding them into injective coefficients and passing to the quotient lowers the degree of a class by
one. This needs the same cochain complex with any functor to abelian groups as coefficients. The
decompositions of test algebras single out the coefficients whose elements are determined by their
images on the pieces.

% before: gal-shifting-cohomology
These groups form a cohomology theory in the coefficients: short exact sequences of coefficient systems
give long exact sequences compatible with pullback, and injective coefficient systems have no
cohomology in positive degrees. In degree $0$ a class is a compatible family of elements, on which local
vanishing and averaging can be checked by hand.

% before: gal-shifting-embedding
The induction on the degree needs an injective coefficient system that still carries decompositions to
products, with a quotient that does too. The coefficient systems that carry decompositions to products
are the abelian sheaves on a site whose coverings are the decompositions, and injective sheaves give
the embedding. The coefficients $A_{\mathrm{cris}}(A)/\ell^\nu$ are among them, since Witt vectors and
divided powers are formed piece by piece.

% before: gal-crystalline-test-cohomology
The cohomology of a functor of points sees only the points, so a weighted sum of maps of points that
cancels locally acts by zero. When a finite group permutes the points over a base, averaging over the
group inverts pullback up to its order. These two properties turn identities between points into
identities between classes.

% before: gal-realization-split-diagrams
The comparison maps go from \'etale cohomology to the cohomology of categories of test algebras. Over a
test algebra $A$, every \'etale cover of $\operatorname{Spec}A^\sharp[1/\ell]$ has a section, so an
\'etale class pulled back there is described by locally constant functions. A diagram of such schemes
indexed by a category carries a coefficient system of these functions, and a Galois element acts on test
algebras by twisting their structure maps.

% before: gal-realization-edge-map
On a scheme whose \'etale covers all have sections, global sections keep an injective resolution of
$\mathbf Z/q$ exact. The rows of the resulting double complex are exact, and the usual staircase carries an
\'etale class of the target to a class with values in locally constant functions. Its naturality and its
compatibility with reduction give the corresponding properties of the comparison maps.

% before: gal-realization-edge-assembly
The staircase in the last proof runs through a double complex whose rows are exact. Its total complex
therefore has the cohomology of the cochains with values in locally constant functions, and the staircase
map is the inverse of this isomorphism composed with the map from \'etale cohomology.

% before: gal-realization-locally-constant
A locally constant function on $\operatorname{Spec}A^\sharp[1/\ell]$ is a sum of multiples of
idempotents, and these idempotents come from $A$. Their Teichm\"uller representatives are fixed by
Frobenius, which is the source of the relation $\phi\circ\eta_\nu=\eta_\nu$.

% before: gal-realization-eta
A point of an $\mathcal O_K$-scheme with values in the untilt of a test algebra gives, after inverting
$\ell$, a geometric point of the generic fibre. These points form a split diagram, and composing its
staircase map with the map to $A_{\mathrm{cris}}$ gives the comparison maps.

% before: gal-realization-duality
The injectivity in Theorem~\ref{gal-crystalline-realization}(2) is proved by pairing a class with classes
of complementary degree. This needs Poincar\'e duality on the geometric generic fibre, and a reduction to
schemes whose generic fibres are connected.

% before: gal-realization-period-datum
To pair classes, Alpöge passes through an auxiliary cohomology with coefficients in
$A_{\mathrm{cris}}$ that maps to the cohomology of $Y^\sharp$ with kernel and cokernel killed by powers of
$t$, and that has periods in $B_{\mathrm{dR}}$. The injectivity argument uses only the multiplicativity
and Galois equivariance of these maps, the bounds by powers of $t$, and the nonvanishing period of a
power of a Chern class.

% before: gal-realization-injectivity
A period datum gives every \'etale class periods in $B_{\mathrm{dR}}$, by lifting $t^n\eta(v)$ to the
auxiliary cohomology. These periods are multiplicative and do not all vanish in top degree. Poincar\'e
duality then shows that a relation among the classes $\eta(v_k)$ forces the same relation among the $v_k$.

% before: gal-realization-descent
In the proof of (3), a Galois-invariant combination is moved into $A_{\mathrm{cris}}\otimes L$, where $L$
is the crystalline cohomology of the special fibre, and must be shown to come from $L$. Since
$B_{\mathrm{cris}}$ is $A_{\mathrm{cris}}$ with $t$ inverted, this is Galois descent after inverting $t$
and $\ell$.

% before: gal-realization-invariants
The descent needs the Galois invariants of $B_{\mathrm{cris}}$ and the Galois action on Witt vectors of
the residue field. Both come from Fontaine's theory of period rings.

% before: gal-realization-frobenius-datum
The proof of (3) passes from the cohomology of $Y^\sharp$ to an auxiliary cohomology with Frobenius, from
there to $A_{\mathrm{cris}}$ tensored with the crystalline cohomology of the special fibre, and back to
points with values in the rings $A/\ell^\flat$. Each passage is an isomorphism up to bounded powers of
$t$, of $\ell$ and of Frobenius, and these bounded losses are what the argument has to control.

% before: gal-realization-twist
A Galois element acts on the cohomology of test algebras by twisting their structure maps, and on
\'etale cohomology through the geometric generic fibre. To compare the two actions, the untilt and the
crystalline ring of a twisted test algebra are identified with those of the original one, compatibly with
locally constant functions and with the geometric points of the comparison maps.

% before: gal-realization-cup-staircase
The comparison maps are multiplicative because the staircase respects products. A cup product of
\'etale classes is represented by a product of representatives, and multiplying a staircase by a cocycle,
on either side, gives a staircase for the product.

% before: gal-realization-period-prism
Most identities below compare two functors into a category of thickenings that differ by a natural family
of maps. Such functors induce the same maps on the cohomology of categories, by the prism construction for
nerves.

% before: gal-realization-period-functors
The comparison bounds below are proved by gluing local computations over a Zariski cover, so the auxiliary
cohomology is set up for functors on rings that need not come from schemes, such as the subfunctors of the members of
a cover. Such a functor only has to say over which map from the base ring each of its points lies. Its points with
values in untilts of test algebras carry the twisted Galois action of the scheme case.

% before: gal-realization-period-thickenings
Alpöge's auxiliary cohomology with coefficients in $A_{\mathrm{cris}}$ is the cohomology of the category of
divided power thickenings of points, in the manner of crystalline cohomology. A test algebra $A$ gives the
thickening $A_{\mathrm{cris}}(A)/\ell^\nu$ of $A^\sharp/\ell^\nu$, and pulling back along these thickenings maps
the auxiliary cohomology to the cohomology of $X^\sharp$. There are two versions: thickenings of points with
values in $\mathcal O_C/\ell^\nu$-algebras, and, with a Frobenius, of points with values in $\mathcal
O_C/\ell$-algebras.

% before: gal-realization-period-compatible
The maps between the auxiliary cohomologies and the maps $\beta_\nu$ must respect reductions, Galois
actions, Frobenius and products, so that they pass to the limit over $\nu$ and give the structures required
of period and Frobenius data. Each identity compares two functors related by a natural family of maps of
thickenings. The divided power filtration is respected as well, and multiplication by an element of the divided power
ideal raises it by one.

% before: gal-realization-period-limit
The bounds on kernels and images are proved one level $\nu$ at a time, while period and Frobenius data
use inverse limits over $\nu$. Passing to the limit squares the bound.

% before: gal-realization-period-t-mu
The bounds at each level are powers of $\mu=[\varepsilon]-1$, while period and Frobenius data are bounded
by powers of $t$. In $A_{\mathrm{cris}}$ the two elements differ by a unit.

% before: gal-realization-period-de-rham
The periods of a period datum are pairings with de Rham cohomology of the generic fibre. Alpöge forms it
as the cohomology of nilpotent divided power thickenings of points, the analogue over a field of the
thickenings above.

% before: gal-realization-period-chern
Condition (4) of a period datum compares the first Chern class of a line bundle in three cohomologies.
Alpöge describes a line bundle by units on a smooth cover, and a first Chern class by the cocycle of
truncated logarithms of these units on every family of thickenings or of points over untilts.

% before: gal-realization-frobenius-special-fibre
For a Frobenius datum, the auxiliary cohomology with Frobenius is compared with the crystalline cohomology
of the special fibre over $W(\mathbf F_q)$. The comparison goes through a power of Frobenius: after the
$\ell^N$-th power map, a point with values in an $\mathcal O_C/\ell$-algebra becomes a point of the special
fibre.

% before: gal-realization-frobenius-fixed
Classes from the special fibre reach test algebras only through the rings $A/\ell^\flat$, on which the
$q$-power Frobenius acts. Pulling them back gives item (2) of a Frobenius datum.

% before: gal-realization-ingredients-cech-line
Item (2) of Lemma~\ref{gal-realization-period-ingredients} needs a line datum whose first Chern classes can be
computed. Alpöge builds line data from cocycles of units on an open cover: the disjoint union of the members of
the cover maps to the scheme, and the transition functions give the units.

% before: gal-realization-ingredients-tautological
On a projective scheme the cocycle of the coordinate ratios $x_k/x_j$ gives a tautological line datum. It is
\'etale and split on the members of a cover, and it is stable under base change, which is what the statements on
Chern classes use.

% before: gal-realization-ingredients-atlas
The comparison bounds are proved locally, on functors represented by rings that are \'etale over a polynomial ring
through units, and then glued. A Zariski atlas records such a cover of a point functor, with the dimension that
controls the bounds.

% before: gal-realization-ingredients-atlas-scheme
The bounds will be applied to smooth projective schemes over $\mathcal O_K$. Smooth, separated and quasi-compact
schemes have Zariski atlases: \'etale coordinates can be shifted to units near each point, and finitely many such charts
cover.

% before: gal-realization-ingredients-bounds
Item (1) of Lemma~\ref{gal-realization-period-ingredients} bounds the kernels and cokernels of the maps
$\beta_\nu$ by one power of $\mu$ for all $\nu$. Its published counterpart compares $A_{\mathrm{inf}}$- and prismatic
cohomology with crystalline cohomology up to bounded powers of $\mu$. Alpöge proves it for the cohomologies of this
section, with an exponent that depends only on the degree and the dimension of an atlas.

% before: gal-realization-ingredients-crystalline-chern
On thickenings, a line datum that is split on the members of an open cover has exactly one first Chern class at
each level. It is the crystalline first Chern class of Berthelot and Illusie.

% before: gal-realization-ingredients-chern-reduction
To form an element of $H^2_\infty$, the first Chern classes of the levels must be compatible under reduction.
Truncating the logarithm at $\ell\nu$ terms loses nothing, because the omitted factorials vanish in
$A_{\mathrm{cris}}/\ell^\nu$.

% before: gal-realization-ingredients-sharp-units
Over untilts, a lift of a line datum needs units of a test algebra whose Teichm\"uller representatives have
given images in the untilt. Every unit of the untilt has this form, because the untilt is absolutely integrally
closed and complete.

% before: gal-realization-ingredients-roots
Item (2)(e) of Lemma~\ref{gal-realization-period-ingredients} compares the first Chern class on thickenings with a
class from Kummer theory, which works with $q$-th roots of units. Two $q$-th roots of a unit differ by a $q$-th root
of unity, and over a scheme with a primitive root of unity such an element is locally a power of it.

% before: gal-realization-ingredients-kummer-maps
Kummer theory produces classes with coefficients in $\mu_q$ and its tensor powers. To write down cochains, each unit
must give a section of an injective resolution of $\mu_q$, and multiplying the unit by a primitive root of unity must
add that root to the section.

% before: gal-realization-ingredients-twisted-edge
The comparison maps of Definition~\ref{gal-realization-eta} use coefficients $\mathbf Z/q$. On the schemes of a split
diagram that carry compatible primitive roots of unity, the twisted sheaves become constant, so the same construction
applies to them, with the same naturality and cup products.

% before: gal-realization-ingredients-kummer-class
A line datum defines units on pairs of points over the scheme, and their $q$-th roots define a class in
$H^2(X,\mu_q)$ through the cohomology of the fibre powers of $f$. The proof of item (2)(e) uses this class through a
representative and two families of cochains that record the choices of roots.

% before: gal-realization-ingredients-kummer-cocycle
On a split diagram the roots can be chosen on each scheme, and their failure to form a cocycle is a locally constant
function with values in $\mathbf Z/q$. This function is a cocycle that represents the twisted comparison of the class
of Lemma~\ref{gal-realization-ingredients-kummer-class}.

% before: gal-realization-ingredients-logarithm
Over test algebras, a first Chern class is written with the truncated logarithm of the Teichm\"uller representatives
of ratios of units. It is additive and compatible with idempotents, and it sends $[\varepsilon]$ to the period $t$.

% before: gal-realization-ingredients-log-roots
The ratios of a lift along test algebras are $q$-th powers of units whose images in the untilt are locally powers of
a fixed root of unity. Their logarithms are $t$ times the corresponding locally constant functions, which connects the
Kummer cocycle with the first Chern class on thickenings.

% before: gal-realization-kappa-approximations
The maps $\kappa$ of Lemma~\ref{gal-realization-period-ingredients}(3) take values in $B_{\mathrm{dR}}^+$, and an
element of $B_{\mathrm{dR}}^+$ is determined by its images in the quotients by the powers of $t$. Alpöge builds $\kappa$
one quotient at a time. Each quotient comes with a ring $A$ with divided powers that receives $A_{\mathrm{cris}}$ and
$\mathcal O_K$ and embeds into it.

% before: gal-realization-kappa-lattice
Alpöge builds the ring $A$ of each approximation inside a truncation of $B_{\mathrm{dR}}^+$, generated by the image of
$W(\mathcal O_C^\flat)$ and the divided powers of $\xi/\ell^M$. The division by a power of $\ell$ leaves room for the
elements of $\mathcal O_K$. The properties of $A$ follow from a few conditions on the truncation, which make sense for
any ring map into a $\mathbf Q$-algebra.

% before: gal-realization-kappa-lattice-ring
Under these conditions the lattice ring has the divided powers and the properties asked for in an approximation. They
come from writing its elements as sums of multiples of the $e_i$, and from the isomorphism of $A/I$ with $W/\xi W$.

% before: gal-realization-kappa-lattice-maps
The Galois actions and the transition maps of an approximation system are restrictions of ring maps between
truncations. Such a map preserves lattice rings and their divided powers when it carries $\xi$ to a multiple of $\xi$.

% before: gal-realization-kappa-newton
The elements of $\mathcal O_K$ are put into a lattice ring by Newton's method. A root of the minimal polynomial is found
in the lattice ring, and it is the given element because roots near a simple approximate root are unique.

% before: gal-realization-kappa-fontaine
The truncations are quotients of Fontaine's ring $B_{\mathrm{dR}}^+$. The construction uses the following properties of
$W(\mathcal O_C^\flat)$, $B_{\mathrm{dR}}^+$ and $A_{\mathrm{cris}}$.

% before: gal-realization-kappa-truncations
The quotients of $B_{\mathrm{dR}}^+$ by the powers of $t$ satisfy the conditions of the matching height. For one $M$,
independent of the height, their lattice rings contain the images of $A_{\mathrm{cris}}$ and of $\mathcal O_K$, and the
Galois action and the projections preserve them.

% before: gal-realization-kappa-approximation-system
The construction of $\kappa$ needs approximations of every height, compatible with each other, whose images determine
elements of $B_{\mathrm{dR}}^+$. Such a system exists for every finite extension $K$ of $\mathbf Q_\ell$.

% before: gal-realization-kappa-truncated-periods
In each approximation $B_m$, a class of the cohomology of thickenings and a linear form on de Rham cohomology have a
value, a truncated period. The relation $\rho_m$ records these values, with the linearity, Galois, cup product and
Chern class properties that $\kappa$ must have.

% before: gal-realization-kappa-truncated-periods-exist-diagram
A truncated period is read off from a diagram of cohomology groups: the levels of the cohomology of thickenings, the
cohomology of thickenings over the rings $A_m$ and $B_m$ of an approximation, and de Rham cohomology. The construction
uses only a few properties of such a diagram, so the diagram is defined abstractly first. A period of a family of classes
is a function on linear forms whose combination with a basis of de Rham cohomology is the image of a lift of the family,
up to a power of $\ell$.

% before: gal-realization-kappa-truncated-periods-exist-periods
In an effective diagram every compatible family has exactly one period, since the cohomology over $B$ is free on the
image of a basis of de Rham cohomology. The linearity, the behaviour under maps of diagrams and the product formulas of
truncated periods follow from the corresponding properties of the diagram.

% before: gal-realization-kappa-truncated-periods-exist-reduction
The cohomology of thickenings of level $\nu$ is formed over $A_\nu=A_{\mathrm{cris}}/\ell^\nu$, so thickenings over an
approximation ring need maps from $A_\nu$. The map from $A_{\mathrm{cris}}$ provides them modulo $\ell^\nu$, and the
conditions on an approximation make them compatible with the divided powers and with $\mathcal O_K$.

% before: gal-realization-kappa-truncated-periods-exist-categories
For each approximation and each smooth projective $\mathcal O_K$-scheme $Y$, the diagram is built from the thickenings of
$Y$ over $A_m$ and of $Y_K$ over $B_m$ whose divided power filtrations vanish from a step $s$ on. Reducing a thickening
over $A_m$ modulo $\ell^\nu$ gives a thickening of level $\nu$; the divided powers pass to the quotient because the
filtration has no $\ell$-torsion.

% before: gal-realization-kappa-truncated-periods-exist-base-change-zariski
The categories of truncated thickenings used here also carry the condition on $\ell$-torsion of the divided power
filtrations. Their cohomology can be computed locally on the scheme, as for the categories without this condition,
because the condition passes to localizations.

% before: gal-realization-kappa-truncated-periods-exist-base-change-splitting
The freeness over $B_m$ and the finiteness over $A_m$ both come from comparing thickenings over a divided power algebra
with thickenings over its base ring. The comparison uses a splitting of the divided power filtration of the algebra into
flat modules over the base.

% before: gal-realization-kappa-truncated-periods-exist-base-change-splitting-field
Over a field such a splitting is a choice of complements. Every element of the algebra then decomposes uniquely along
the splitting, compatibly with the filtration.

% before: gal-realization-kappa-truncated-periods-exist-base-change-chart
The local comparison runs on divided power charts of the truncated categories. Over a field of characteristic $0$ they
were built for de Rham cohomology; over $A_m$ the integers need not be invertible, and the divided powers on a chart come
from the absence of torsion instead.

% before: gal-realization-kappa-truncated-periods-exist-base-change-tensor
With a splitting, the cohomology of the truncated thickenings over the algebra, with coefficients in a quotient of the
divided power filtration, is the sum of the tensor products of the splitting modules with the corresponding cohomology
over the base. The proof reduces to a coordinate open by Zariski descent and \v{C}ech gluing, and there to the charts.

% before: gal-realization-kappa-truncated-periods-exist-base-change-filtration
Over a field of characteristic $0$, the cohomology of the truncated thickenings with coefficients in a quotient of the
filtration is de Rham cohomology once the step exceeds the dimensions. A de Rham class maps to $0$ there if and only if
it is the class of a cocycle with values in the corresponding step of the filtration.

% before: gal-realization-kappa-truncated-periods-exist-base-change
For the diagram to be effective, the cohomology of the truncated thickenings over $B_m$ must be free over $B_m$ on the
image of a basis of de Rham cohomology. This holds once the step $s$ exceeds the dimensions of a coordinate cover by the
height of the approximation, and the comparison respects the divided power filtrations.

% before: gal-realization-kappa-truncated-periods-exist-lifting-limits
Lifting a compatible family of classes modulo the powers of $\ell$ to a class over $A_m$ is a statement about a complex
without $\ell$-torsion and its reductions. It holds up to a fixed power of $\ell$ once the cohomology has bounded
$\ell$-power torsion, its $\ell$-divisible elements are killed by a fixed power of $\ell$, and its Cauchy sequences for
$\ell$ have limits up to such a power.

% before: gal-realization-kappa-truncated-periods-exist-lifting-modules
These three properties hold for modules over $A_m$ that are finitely generated up to a power of $\ell$. The conditions
on $A_m$ that enter are that $I_m$ is a finitely generated nilpotent prime ideal, that a power of $\ell$ lies in
$xA_m+I_m$ for every $x\notin I_m$, and that $A_m$ is complete for $\ell$.

% before: gal-realization-kappa-truncated-periods-exist-lifting-splitting
The finite generation over $A_m$ comes from the comparison with thickenings over $\mathcal O_K$, which needs a graded
splitting of $A_m$ over $\mathcal O_K$. The graded pieces of the filtration of $A_m$ are torsion-free over the discrete
valuation ring $\mathcal O_K$, and a section of each is built from a free submodule by completeness.

% before: gal-realization-kappa-truncated-periods-exist-lifting-finite-local
The finiteness over $\mathcal O_K$ needs the local computation of the graded cohomology, which was done over a field of
characteristic $0$ for de Rham cohomology. Over a base without $\ell$-torsion the same computation runs on the charts
built from the absence of torsion.

% before: gal-realization-kappa-truncated-periods-exist-lifting-finite
The cohomology of the truncated thickenings over $\mathcal O_K$ is finitely generated over $\mathcal O_K$ for a proper
scheme. This is the finiteness of de Rham cohomology over a field, with $\mathcal O_K$ in place of the field.

% before: gal-realization-kappa-truncated-periods-exist-lifting-generation
Combining the comparison with the finiteness over $\mathcal O_K$ shows that the cohomology over $A_m$ is finitely
generated over $A_m$. The steps of the filtration of $A_m$ are finitely generated ideals, which bounds the number of
generators.

% before: gal-realization-kappa-truncated-periods-exist-lifting
Effectiveness also asks that compatible families of classes modulo the powers of $\ell$ come from classes over the
approximation ring $A_m$, up to a power of $\ell$ independent of the family. The conditions on $A_m$ in an approximation,
which make $A_m/I_m$ behave like a valuation ring complete for $\ell$, enter here.

% before: gal-realization-kappa-truncated-periods-exist-effective
With these two inputs and the finiteness of de Rham cohomology, the diagram of an approximation is effective once the
step $s$ is large. The remaining conditions say that the reduction maps are compatible with the levels.

% before: gal-realization-kappa-truncated-periods-exist-maps
The properties of truncated periods under a change of the step, of the approximation, of the Galois element or of the
scheme come from maps between the diagrams. Each square of such a map compares two functors on thickenings, which agree
or are related by identity maps of rings.

% before: gal-realization-kappa-truncated-periods-exist-chern-log
The class over $A_m$ that the Chern class property needs is built from the logarithm on the divided power ideals of the
thickenings. On a thickening over $A_m$ the integers need not be invertible, so the additivity of the truncated
logarithm is proved from the divided power identities alone, once its terms of high order vanish.

% before: gal-realization-kappa-truncated-periods-exist-chern-cocycle
With the additive logarithm, a lift of a line datum along a family of thickenings gives a $2$-cocycle, as in de Rham
cohomology. The argument used there over a field applies to thickenings over any base, and two lifts give cohomologous
cocycles.

% before: gal-realization-kappa-truncated-periods-exist-chern-class
These cocycles define first Chern classes in the cohomology of a category of truncated thickenings, by the condition used
in de Rham cohomology. The condition is tested only on families of thickenings over opens where the line datum has a
section, so that lifts exist.

% before: gal-realization-kappa-truncated-periods-exist-chern-unique
When such opens cover the scheme, the cohomology of the category is computed by the thickenings over them, by Zariski
descent. A first Chern class then exists and is unique.

% before: gal-realization-kappa-truncated-periods-exist-chern-reduction-functor
To compare the class over $A_m$ with a first Chern family, the thickenings over $A_m$ are reduced modulo $\ell^\nu$ to
thickenings of level $\nu$. The reduction works over any base with a ring map to $A_m/\ell^\nu A_m$ that carries its
divided power ideal into the image of $I_m$, compatibly with the divided powers and with the structure maps.

% before: gal-realization-kappa-truncated-periods-exist-chern-reduction
Along such a reduction, a first Chern class over $A$ reduces modulo $\ell^n$ to a first Chern class over the base, as the
logarithm commutes with the reduction and its terms of high order vanish modulo $\ell^n$. At level $\nu$ this compares the
class over $A_m$ with the components of a first Chern family.

% before: gal-realization-kappa-truncated-periods-exist-chern-pullback
A divided power map of bases with a compatible morphism of schemes carries a lift of a line datum to a lift of its base
change with the same cocycle. So first Chern classes correspond under the induced functor, and for $A_m\to B_m$ and
$Y_K\to Y$ the first Chern class over $A_m$ maps to the first Chern class over $B_m$.

% before: gal-realization-kappa-truncated-periods-exist-chern-de-rham
It remains to identify the class over $B_m$ with the image of a first Chern class in de Rham cohomology. A thickening over
$B_m$ is in particular a nilpotent thickening over $K$, so the defining property of the de Rham class applies to it.

% before: gal-realization-kappa-truncated-periods-exist-chern
The Chern class property compares a first Chern family of a tautological line datum with a first Chern class in de Rham
cohomology. Both come from one class in the cohomology of the thickenings over $A_m$, which reduces to the first and
maps to the image of the second.

% before: gal-realization-kappa-truncated-periods-exist-filtration
The last property, on filtrations, is not used in this book but is part of the definition. A filtered class has a lift
whose image vanishes modulo the corresponding step, and the filtered form of the freeness over $B_m$ bounds its periods.

% before: gal-realization-kappa-truncated-periods-exist
Truncated periods exist for every approximation system. With the uniqueness in the approximation system, they define
$\kappa(x)(\psi)$ as the element of $B_{\mathrm{dR}}^+$ with the given images in all the $B_m$.

% before: gal-realization-kappa-de-rham-truncated
The finiteness of de Rham cohomology is proved by the method used at level $1$ for the cohomology of thickenings over
$W_1$. Over a field the divided power filtration of a thickening need not stop, so the thickenings are truncated: those
whose filtration vanishes from a given step on carry the same cohomology once the step exceeds the dimension.

% before: gal-realization-kappa-de-rham-subideal
Truncating a thickening means dividing it by a step of its divided power filtration, and the quotient keeps divided
powers when that step is stable under them. The divided powers of an element of a step lie in a deeper step, which
gives this stability for every element of the step.

% before: gal-realization-kappa-de-rham-poly
The divided power polynomial structure of the thickenings, which drives the induction on the dimension of a torsor
datum, must stay inside the truncated thickenings. Dividing the divided power polynomial algebra by the same step of its
filtration achieves this, and the criteria for the filtration survive up to that step.

% before: gal-realization-kappa-de-rham-chart
Coordinates on an affine open give a divided power chart over any base in which every integer divides a power of $\ell$
and $\ell$ is a nonzerodivisor modulo each step of the divided power filtration, for instance one in which the integers
are invertible. Its ring is the tensor product of the base with the ring of functions, its divided powers are the
quotients $x^n/n!$, and its divided power filtration is generated by that of the base, so the chart is truncated when
the base is.

% before: gal-realization-kappa-de-rham-local
Over an affine open with coordinates, the chart built from the coordinates over the field is the ring of functions, with
zero divided power ideal, and it lies among the truncated thickenings. As at level $1$, the graded piece of weight $r$ then has cohomology only in degree $r$, where it is
free on the forms in the coordinates, and none at all when $r$ exceeds the dimension.

% before: gal-realization-kappa-de-rham-bundle
The freeness on affine opens gives, for a family of opens with coordinates and affine intersections, the gluing data of a
vector bundle for the \v{C}ech complex of the graded pieces in their own degree. The argument is the one used at level
$1$.

% before: gal-realization-kappa-de-rham-zariski
The truncated thickenings are stable under localization, since localization commutes with the divided power filtration.
So the comparison between thickenings over two families of opens with the same union holds for them too.

% before: gal-realization-kappa-de-rham-affine-basic
The vanishing on an affine open is proved on a cover of the open by basic opens with coordinates. Such a cover comes from
the local structure of smooth algebras: a smooth algebra is standard smooth after inverting each of finitely many
elements, and a standard smooth algebra is \'etale over a polynomial algebra. Intersections of basic opens are basic opens,
so they are affine.

% before: gal-realization-kappa-de-rham-affine
On an affine open the vector bundle has no higher cohomology. Together with the local vanishing outside degree $r$, this
shows that over an affine open the graded piece of weight $r$ has cohomology only in degree $r$, with or without
coordinates.

% before: gal-realization-kappa-de-rham-tower
To pass from all nilpotent thickenings to the truncated ones, the cohomology of a deep step of the filtration must
vanish. The filtration vanishes at each thickening from some step on, which may depend on the thickening, so the
vanishing of the graded pieces is assembled by a convergent sum of cochains.

% before: gal-realization-kappa-de-rham-restriction
Dividing a thickening by a step of its divided power filtration is left adjoint to the inclusion of the truncated
thickenings. It does not change the quotients of the filtration below that step, so their cohomology is the same on all
thickenings and on the truncated ones.

% before: gal-realization-kappa-de-rham-truncation
These results identify de Rham cohomology with the cohomology of the truncated thickenings over a cover, once the
truncation step exceeds the dimensions of the coordinates. The deep steps of the filtration have no cohomology, because
their graded pieces vanish locally.

% before: gal-realization-kappa-de-rham-log
A first Chern class in de Rham cohomology is represented by the truncated logarithm of the transition units of a line
datum. On a thickening whose divided power filtration vanishes at the truncation step, this logarithm turns products of
units congruent to $1$ into sums. When the integers are invertible it is a polynomial whose derivative is a truncated
geometric series, which gives the identity.

% before: gal-realization-kappa-de-rham-lift
The cochain of a first Chern class is formed from a lift of the line datum: a point of its total space over each
thickening and a unit for each map. Over an open where the datum has a section the points come from the section, and
the units lift because the divided power ideal of a thickening is nilpotent.

% before: gal-realization-kappa-de-rham-cocycle
The truncated logarithm of the units of a lift is a cocycle, because the units satisfy a multiplicative cocycle identity
and the logarithm is additive. Two lifts change it by the coboundary of the logarithm of the ratio of their units, so its
class depends only on the line datum.

% before: gal-realization-kappa-de-rham-cone
To show that the top power of the Chern class is not zero, the scheme is covered by the opens where one of finitely many
homogeneous polynomials of the same degree does not vanish. Each polynomial trivializes the corresponding power of the
tautological line datum on its open, and the ratios of the polynomials are the transition functions, whose logarithmic
derivatives are cocycles.

% before: gal-realization-kappa-de-rham-torus
The contradiction at the end of the argument is an identity among functions on an affine open whose coordinates are
ratios of the polynomials. It is written with the algebras of Laurent monomials in these coordinates and with the
Jacobian determinants of functions from the intersections of the opens.

% before: gal-realization-kappa-de-rham-cone-vanishing
The graded cohomology of the categories attached to the intersections of the opens is that of the truncated thickenings
over these intersections, since the standard charts of projective space cover the scheme. On an affine open it vanishes
outside degree $r$, and above the dimension it vanishes in every degree.

% before: gal-realization-kappa-de-rham-setup-normalization
The polynomials $y_k$ come from Noether normalization of the homogeneous coordinate ring of the scheme. Taking them
homogeneous of one degree makes their ratios functions on the intersections of the opens where they do not vanish.

% before: gal-realization-kappa-de-rham-setup-cone-relation
Integrality of a ratio of forms over the Laurent monomials comes from an integral equation in the homogeneous coordinate
ring. Taking a homogeneous component of such an equation makes each coefficient a form of the right degree in the
polynomials $y_k$.

% before: gal-realization-kappa-de-rham-setup-cone-vanish
The ideal of the scheme is described by the vanishing of ratios of forms. A form in the ideal gives a ratio that vanishes
on every open where its denominator does not, and on an integral scheme a ratio that vanishes on one nonempty open comes
from a form in the ideal. The ideal also lies in the homogeneous prime of every point.

% before: gal-realization-kappa-de-rham-setup-cone-monomial
On the intersection of the opens, a monomial in the $y_k$ divided by a power of the product $y_S$ is a Laurent monomial in
the ratios $y_{k+1}/y_0$. Its exponents satisfy the sign conditions that define the algebras of Laurent monomials, and
every such monomial arises this way.

% before: gal-realization-kappa-de-rham-setup-cone
The normalization gives a cover of the scheme by the opens where one of the $y_k$ does not vanish, with affine
intersections. On each intersection the ratios $y_{k+1}/y_0$ are units, and every function is integral over the Laurent
monomials in them that are regular there.

% before: gal-realization-kappa-de-rham-setup-generic
Four general facts are used to shrink the intersection of all the opens to an open with coordinates. In characteristic
$0$ a domain algebraic over a polynomial algebra is \'etale over it after inverting one element, and the number of
coordinates is then the relative dimension.

% before: gal-realization-kappa-de-rham-setup
The polynomials are chosen so that their opens cover the scheme with affine intersections and their ratios are
coordinates on an affine open. The rings of the intersections then embed in the ring of that open, finite over algebras
of Laurent monomials in the coordinates.

% before: gal-realization-kappa-de-rham-residue-torus
The final contradiction is proved by taking traces down to the field of rational functions in the units and then the
coefficient of $(t_0\cdots t_n)^{-1}$. This needs the algebras of Laurent monomials on each face to be localizations of
polynomial algebras in suitable coordinates, and the monomials of each face to avoid the exponent $(-1,\ldots,-1)$.

% before: gal-realization-kappa-de-rham-residue-functional
The coefficient of $(t_0\cdots t_n)^{-1}$ is a linear form on the Laurent monomials. It extends to the whole ring, and it
kills derivatives of Laurent polynomials, as a derivative of a monomial never has the exponent $(-1,\ldots,-1)$.

% before: gal-realization-kappa-de-rham-residue-finite
For the trace to be defined, the field of fractions of the ring has to be finite over the field generated by the units.
This holds because the algebra of the whole intersection is integral over the Laurent polynomials and generates the
ring by fractions.

% before: gal-realization-kappa-de-rham-residue-chart
On a face, the given algebra need not generate the ring by fractions. Enlarging it by suitable multiples of the
generators of the algebra of the whole intersection repairs this while keeping it integral over the Laurent monomials of
the face.

% before: gal-realization-kappa-de-rham-residue-jacobian-poly
The trace of a Jacobian determinant is shown to be integral one prime at a time, over the discrete valuation ring at a
prime element of the polynomial algebra. There the derivative of an integral element is controlled through a polynomial
that vanishes at it modulo a prime but whose derivative does not. Such a polynomial exists once the integers are
invertible.

% before: gal-realization-kappa-de-rham-residue-jacobian-valuation
A derivation that preserves the discrete valuation ring and takes its uniformizer into the uniformizer times integers
does not create poles on the integral closure. If moreover it takes the ring into an ideal inside a prime, its values on
the integral closure lie in that prime. The bound is improved one step at a time from a crude bound given by a common
denominator.

% before: gal-realization-kappa-de-rham-residue-jacobian-nilpotent
An integral element some power of which is divisible by the uniformizer acts on the integral closure, a free module over
the valuation ring, by a matrix that is nilpotent modulo the uniformizer. So its trace is divisible by the uniformizer.

% before: gal-realization-kappa-de-rham-residue-jacobian-dvr
Over a discrete valuation ring, the Jacobian determinant is multiplied by the uniformizer after row operations that make
one derivation take the ring into the uniformizer times the ring and the others kill the uniformizer. The product is then
in every prime above the uniformizer, so a power of it is divisible by the uniformizer, and its trace is too.

% before: gal-realization-kappa-de-rham-residue-jacobian
The trace of a Jacobian determinant formed with derivations dual to the coordinates of a localized polynomial algebra
lands in that algebra when the entries are integral over it. This is what keeps the traces of the Jacobian terms
inside the Laurent monomials of each face.

% before: gal-realization-kappa-de-rham-residue-trace
The Jacobian terms of the identity are minors, so the coordinates missing from a minor are added as entries to reach a
full determinant. The fields of fractions are identified so that the previous lemma applies on each face.

% before: gal-realization-kappa-de-rham-residue-derivation
The divergence terms $D_k(\eta_k)$ are handled by moving the derivation outside the trace. This is possible when the
derivation preserves the base field.

% before: gal-realization-kappa-de-rham-residue
The final contradiction is a statement about a domain with algebraically independent units and dual derivations. The
inverse of the product of the units is not a sum of Jacobian terms from the faces and of derivatives of Jacobian terms
from the whole intersection.

% before: gal-realization-kappa-de-rham-local-exact-charts
A thickening over a standard chart of projective space gives a ring map from the degree-zero fractions of that chart to
the reduction of its ring. The local primitives below are built from the values of these maps on the fractions
$y_k/x_j^e$, so their behaviour under maps of thickenings and under a change of chart is needed first.

% before: gal-realization-kappa-de-rham-local-exact
On the open where $y_k$ does not vanish, $y_k$ trivializes the $e$-th power of the tautological line datum, so $e$
times the Chern cocycle should be a coboundary on the corresponding category. The primitives on two such categories
differ by the logarithmic derivative of the ratio of the polynomials.

% before: gal-realization-kappa-de-rham-cup
Cup products of cocycles with values in the first step of the filtration take values in the corresponding deeper step,
so they define classes in the graded cohomology. The logarithmic derivatives of the transition functions are such
cocycles.

% before: gal-realization-kappa-de-rham-cech-stair
The relation among the classes on the intersections comes out of a \v{C}ech double complex with a filtration. If the
filtration can be raised by one step on a column whenever the vertical coboundary allows it, an element of the corner
that is a total coboundary splits into a horizontal and a vertical coboundary, each one step down the filtration.

% before: gal-realization-kappa-de-rham-cech-zigzag
The element of the corner is reached from the top cup power of the cocycle by a zigzag. Each step multiplies by the
difference of two local primitives and moves one place down the diagonal.

% before: gal-realization-kappa-de-rham-cech
If the top power of a class vanishes while the class is a coboundary on each piece of a cover, the cup product of the
differences of the local primitives is constrained on the intersection of all the pieces. It becomes a sum of classes
from the faces and of the class of a coboundary of a cochain one step down the filtration, provided the graded
cohomology of the intersections vanishes outside the expected degrees.

% before: gal-realization-kappa-de-rham-transfer-lifts
To compare the relation on the categories $\mathcal W_S$ with classes on the truncated thickenings over the opens $U_S$,
the logarithmic derivatives have to be formed with lifts chosen on the larger categories. Their graded classes do not
depend on these choices.

% before: gal-realization-kappa-de-rham-transfer-extend
The cochain $f$ of the relation lives on $\mathcal W_\top$. When restriction is surjective in graded cohomology one degree
lower, it can be replaced by a cochain on the larger category whose coboundary has the restricted class.

% before: gal-realization-kappa-de-rham-transfer-local-near
Over a chart, cocycles are compared with the forms $\omega_I$ of a graded constant torsor datum. When each of several
$1$-cochains agrees with a combination of the forms $\omega_{(i)}$ up to a coboundary and higher filtration, their cup
product agrees in the same way with the corresponding combination of the forms $\omega_I$, and has the same graded class.

% before: gal-realization-kappa-de-rham-transfer-local-alternating
The forms $\omega_I$ are cup products of the $1$-cocycles $a_k$, and in graded cohomology these products alternate like a
wedge product. So a sum of forms indexed by arbitrary sequences has the class of a sum indexed by strictly increasing
sequences, as in Lemma~\ref{gal-realization-thickening-level-one-graded}(2), with signs of permutations as coefficients.

% before: gal-realization-kappa-de-rham-transfer-local-connecting
A first-order Taylor expansion of the maps from the chart computes the coboundary of a form $\omega_I$ modulo the next
step of the filtration. It gives the graded class of the coboundary of a cochain from the coefficients of its own class,
as a de Rham differential would.

% before: gal-realization-kappa-de-rham-transfer-local-dlog
The same expansion, applied to a unit, gives its logarithmic derivative in graded cohomology. Up to a coboundary and
higher filtration, the ratio of the lifts of the unit along a map is a combination of the forms $\omega_{(i)}$ whose
coefficients are the logarithmic derivatives of the unit along the coordinates.

% before: gal-realization-kappa-de-rham-transfer-local-taylor
The Taylor expansions rest on an algebraic fact. Two algebra maps that agree to first order on an algebra formally
unramified over a polynomial algebra differ, to second order, by the derivatives along the coordinates.

% before: gal-realization-kappa-de-rham-transfer-local-chart
The chart of Lemma~\ref{gal-realization-kappa-de-rham-local} is needed on several affine opens with coordinates at once.
On each, graded classes of degree $m<s$ are given by unique functions, and the maps from the chart to the thickenings
satisfy a first-order Taylor expansion.

% before: gal-realization-kappa-de-rham-transfer-local-compare
When an affine open with coordinates lies in another, a graded class on the larger one restricts to the smaller one, and
the functions that represent it in the two charts are related by Jacobian determinants of the coordinates. This is the
change of variables formula for differential forms of degree $m$.

% before: gal-realization-kappa-de-rham-transfer-local-clear
The change of variables holds on a basic open, where the denominators of the second chart are inverted. An algebraic
lemma removes these denominators at the cost of a power of the function that defines the basic open.

% before: gal-realization-kappa-de-rham-transfer-local-minors
A class that comes from an affine open containing the open with coordinates should have coefficients built from functions
on the larger open. Covering the larger open by basic opens with coordinates, the coefficients lie in the group generated
by the Jacobian minors of such functions.

% before: gal-realization-kappa-de-rham-transfer-local
Over an open with coordinates $t_0,\ldots,t_n$, a top graded class is a function times $dt_0\wedge\cdots\wedge dt_n$. The
cup product of logarithmic derivatives gives the determinant of the logarithmic derivatives, classes from an affine open
give Jacobian terms, and graded coboundaries give divergences.

% before: gal-realization-kappa-de-rham-transfer
Over an affine open with coordinates, classes in the graded cohomology are given by differential forms in the
coordinates. The relation of the previous lemma yields an identity of functions on such an open, between the
determinant of the logarithmic derivatives of the transition functions and Jacobian terms.

% before: gal-realization-kappa-de-rham
The rest of Lemma~\ref{gal-realization-period-ingredients}(3) concerns de Rham cohomology alone: the finiteness of the
groups, which is needed for dual bases, and the nonvanishing of the top power of the Chern class of a tautological line
datum. These are the counterparts, for the cohomology of nilpotent thickenings, of the corresponding facts for
algebraic de Rham cohomology.

% before: gal-realization-frobenius-congruence
The Frobenius factorization of Lemma~\ref{gal-realization-period-ingredients}(4) compares two $\mathcal O_K$-structures
on thickenings: through $\mathcal O_K\subset\mathcal O_C$, and through the residue field and the Witt vectors. A power
of Frobenius carries the first to the second, because $x^{\ell^N}$ is congruent modulo $\ell$ to the Teichm\"uller
representative of the residue of $x$ when $N$ is a large multiple of $f$.

% before: gal-realization-frobenius-teichmuller
Thickenings whose $\mathcal O_K$-structure passes through the residue field sit between the two cohomologies of
Lemma~\ref{gal-realization-period-ingredients}(4). They become thickenings over $W_\nu$ through $W_\nu\to A_\nu$, and
they are related to the thickenings for $I'_\nu$ by the $N$-th powers of Frobenius on the points and on
$A_{\mathrm{cris}}$.

% before: gal-realization-frobenius-teichmuller-maps
The two Frobenius functors compose to the $N$-th power of the Frobenius of $H'^j_\nu(Y)$, and the passage from
thickenings for $I'_\nu$ to thickenings over $W_\nu$ through the Teichm\"uller thickenings is a Frobenius twist. These
identities hold on the categories, so the proofs compare composite functors.

% before: gal-realization-thickening-nilpotent
The base change to $A_{\mathrm{cris}}$, the comparison with flat thickenings and the finiteness at level $1$ all rest on
one theory of the cohomology of thickenings. It works over any ring with divided powers in which a fixed integer is
nilpotent, and it is local on the scheme: the cohomology is computed over affine opens with \'etale coordinates and then
glued.

% before: gal-realization-thickening-over
Opens of the scheme give full subcategories of thickenings. A thickening that lies over an open passes this property to
every thickening it maps to, and it lies over an intersection of opens when it lies over each of them.

% before: gal-realization-thickening-coordinates
Over an affine open with \'etale coordinates, the thickenings have a distinguished object. Its ring is \'etale over a
polynomial ring over the base, its ideal comes from the base, and it represents the points of the open with values in
algebras over the residue ring of the base.

% before: gal-realization-thickening-cover
Smooth schemes are covered by opens with \'etale coordinates. On a separated quasi-compact scheme finitely many of them
suffice, and their finite intersections are affine opens that again carry coordinates.

% before: gal-realization-thickening-chart
A divided power chart maps to every thickening over its open, and these maps are a torsor under the $d$-th power of the
divided power ideal, because the chart is \'etale over a polynomial ring and divided power ideals are nil. Charts with
flat rings exist once the integer $p$ is nilpotent, and charts are compatible with change of the base ring.

% before: gal-realization-thickening-torsor
The local computations use only a few features of the categories of thickenings over an open with coordinates. A chart
gives a torsor datum, and divided power polynomial algebras describe the maps that extend a given map by one element of
the divided power ideal. Both are stated for any category with a functor to rings with divided powers.

% before: gal-realization-thickening-chart-torsor
A chart gives a torsor datum for every category of thickenings over its open that contains it, and restriction of scalars
along a divided power map is a morphism of the corresponding functors. The divided power polynomial algebras give the
polynomial structure, both for all thickenings and for those with flat rings.

% before: gal-realization-thickening-torsor-cut
The comparisons below are proved by induction on the dimension of a torsor datum. The induction passes from a category
to the pairs of an object and a map to it from the object of the datum, and then to the maps along which the first
coordinate of the datum vanishes. These carry a torsor datum of one dimension less.

% before: gal-realization-thickening-torsor-cut-adjoint
The induction needs the cut category to carry the same structures as the original one. The divided power polynomial
structure adjoins a variable for the first coordinate, and this gives a left adjoint to the inclusion of the cut
category into the pointed one.

% before: gal-realization-thickening-torsor-adjoint
Adjoint functors let the cohomology be computed on the pointed and on the cut categories. In dimension $0$ the object of
a torsor datum is initial, and an initial object computes the cohomology directly.

% before: gal-realization-thickening-torsor-exact
Adjoining the variable shifts the degrees of a ladder by one. Differentiation in the variable gives a short exact
sequence that relates the cohomology of the quotients in degrees $r$ and $r-1$ to that of the pointed category.

% before: gal-realization-thickening-torsor-cut-morphism
The comparisons relate two categories with torsor data through a morphism over a ring map whose maps are injective and
which is compatible with the torsor data. Its functor passes to the pointed and cut categories, and it is compatible with
the adjoined variables through a natural transformation.

% before: gal-realization-thickening-torsor-base-change
On a category with a torsor datum and a divided power polynomial structure, the object of the datum controls the
cohomology. So a comparison of cohomology after base change reduces to the same comparison for that object, in as many
degrees as the dimension.

% before: gal-realization-thickening-local-base-change
Over an open with coordinates, the charts over $R$ and over $R'$ are related by extension of scalars. Their torsor data
then give base change for the cohomology of all thickenings over the open.

% before: gal-realization-thickening-cosieves
The thickenings lying over the opens of a cover form full subcategories that are closed under maps out of them. For a
category covered by finitely many such cosieves, base change, bijectivity of maps and vanishing of cohomology pass from
the finite intersections to the whole category.

% before: gal-realization-thickening-zariski-class
A thickening over the union of a family of opens is covered by the localizations of its ring at finitely many elements,
each of which gives a thickening over one of the opens. The abstract form of this situation is a class of objects that
receives maps from every object after localization, with functors on which localization acts as it does on modules.

% before: gal-realization-thickening-zariski-double-up-to
Restriction to a Zariski class is compared through double complexes built from the categories of maps out of a fixed
object. The two facts on double complexes needed for this hold with a multiplier: a map whose kernel and cokernel are
killed by $s$ behaves like an isomorphism, at the cost of powers of $s$.

% before: gal-realization-thickening-zariski-double
Two facts on first quadrant double complexes are used, one for rows that are resolutions and one for maps that are
isomorphisms on the cohomology of the columns. Both are proved by diagram chases in the total complex, and with a
multiplier $s$ each step of the chase costs a factor of $s$.

% before: gal-realization-thickening-zariski-double-filtration
The comparison of total complexes is often proved with the filtration of a double complex by its columns. That proof
uses the hypothesis on the columns in all degrees at once, and the five lemma in place of the chase.

% before: gal-realization-thickening-zariski-fibre
A functor between categories induces isomorphisms on cohomology when, for every object of the target, the category of
maps from it into the image is acyclic with the right value in degree $0$. The proof compares two double complexes built
from these categories of maps.

% before: gal-realization-thickening-zariski-cech
The acyclicity comes from localization. On the nonempty subsets of a finite set, the cochains with values in the
localizations at the products of elements generating the unit ideal form a resolution of the module itself, as the
\v{C}ech complex of a cover by basic opens does.

% before: gal-realization-thickening-zariski-refinement
The categories of maps that arise from a Zariski class reduce to the subsets of the indices of the localizations, through
an adjunction. They are then acyclic by the previous lemma.

% before: gal-realization-thickening-zariski-descent
For such a class the cohomology of the whole category is that of the full subcategory on the class. This is the
counterpart, for these cochain complexes, of the vanishing of higher \v{C}ech cohomology of a module on a cover of an
affine scheme by basic opens.

% before: gal-realization-thickening-zariski-localization
Localizing the ring of a thickening at one element gives a thickening again, since divided powers extend along flat ring
maps. It receives a map from the thickening through which every map that inverts the element factors, and its divided
power filtration is the extension of the filtration of the original ring.

% before: gal-realization-thickening-zariski
Restricting to the thickenings that lie over the opens of a cover does not change cohomology. This reduces statements
about all thickenings of a scheme to a category covered by cosieves, to which
Lemma~\ref{gal-realization-thickening-cosieves} applies.

% before: gal-realization-thickening-base-change
For a scheme with a coordinate cover, the local base change over the finite intersections of the cover glues along the
cosieves of the cover. Zariski descent then identifies the result with the cohomology of all thickenings.

% before: gal-realization-frobenius-base-change
The cohomology of Teichm\"uller thickenings is obtained from the cohomology over $W_\nu$ by extending scalars to
$A_{\mathrm{cris}}$. This is the counterpart of flat base change for crystalline cohomology, and it lets the map
$\nu_\nu$ take values in $A_{\mathrm{cris}}\otimes_WH^j_\nu(Y/W)$.

% before: gal-realization-frobenius-bockstein
Passing from the levels $\nu$ to the limit $L^j(Y)$ needs control of the inverse system of the $H^j_\nu(Y/W)$. It
comes from the long exact sequences attached to $0\to W_a\to W_{a+b}\to W_b\to0$, whose properties are listed here for
any inverse system.

% before: gal-realization-frobenius-artinian-lift
Elements of an inverse limit are built from compatible choices at each level. For Artinian modules a family that is
compatible only modulo given submodules can be corrected to a compatible one, because decreasing chains of submodules
stabilize.

% before: gal-realization-frobenius-bockstein-limits
Bockstein sequences make an inverse system behave like the reductions of a finitely generated module. Finite length
passes from the first level to all levels, the limit is finitely generated, the projections from it are surjective up
to a fixed power of $\ell$, and their kernels are the expected multiples.

% before: gal-realization-frobenius-flat-tower
Bockstein sequences arise from complexes that are flat over the rings $\mathbf Z/\ell^n$ and agree modulo $\ell^n$ up
to quasi-isomorphism. A tower of such complexes is what the cohomology of thickenings over the $W_\nu$ provides.

% before: gal-realization-frobenius-flat-tower-bockstein
For a flat tower, multiplication by powers of $\ell$ gives short exact sequences of complexes at each level. Their long
exact cohomology sequences are the Bockstein sequences, once the levels are identified through the maps of the tower.

% before: gal-realization-thickening-local-flat
Over an open with coordinates, a chart with flat ring lies both among all thickenings and among those with flat rings,
and it gives torsor data for both. The comparison of Lemma~\ref{gal-realization-thickening-torsor-base-change}, with no
change of the base ring, then shows that the flat thickenings have the same cohomology.

% before: gal-realization-thickening-flat-restriction
Gluing over a coordinate cover, as for base change, turns the local comparison into one for the whole scheme. It gives
the bijectivity of the restrictions $\mathrm{res}_m$ in Lemma~\ref{gal-realization-frobenius-flat-thickenings}(5).

% before: gal-realization-thickening-torsor-quotient
Reduction modulo $\ell^n$ carries flat thickenings over $W_{n+1}$ to flat thickenings over $W_n$, and on rings it is an
isomorphism onto a quotient of the original ring. Torsor data can be compared along such functors as well.

% before: gal-realization-thickening-torsor-cut-morphism-quotient
A morphism modulo a family of ideals passes to the pointed and cut categories in the same way as a morphism over a ring
map, with every identity between values now holding modulo the family. Injectivity of its maps is what lets the cut
condition on maps pass from one side to the other.

% before: gal-realization-thickening-torsor-quotient-comparison
As for base change, the cohomology of two categories with torsor data can be compared along a morphism modulo a family of
ideals by comparing the objects of the data. Reduction modulo $\ell^n$ is such a morphism, modulo the multiples of
$\ell^n$.

% before: gal-realization-thickening-chart-reduction
Reducing a chart over $W_{n+1}$ modulo $\ell^n$ gives a chart over $W_n$, because \'etale maps, the ideal generated by
$\ell$ and the residue field $k$ are all compatible with the reduction. The chart and its reduction then give
compatible torsor data at the two levels.

% before: gal-realization-thickening-local-quotient
Over an open with coordinates, a flat chart over $W_{n+1}$ and its reduction give compatible torsor data. The comparison
modulo $\ell^n\mathcal O$ then shows that reducing flat thickenings modulo $\ell^n$ does not change their cohomology over
the open.

% before: gal-realization-thickening-flat-quotient
Gluing over a coordinate cover gives the comparison for the whole scheme. It shows that the maps $\psi_n$ of
Lemma~\ref{gal-realization-frobenius-flat-thickenings} induce isomorphisms on cohomology.

% before: gal-realization-frobenius-flat-witt
The flatness and reduction properties below hold for thickenings over truncated Witt vectors of any perfect field, of
any point functor whose points lie over the residue field. Only the comparison with all thickenings needs the functor
of a smooth projective scheme.

% before: gal-realization-frobenius-flat-thickenings
The thickenings whose rings are flat over $W_m$ have cochain complexes that are flat in this sense, and reducing them
modulo $\ell^n$ links consecutive levels. For a smooth projective scheme these complexes compute the cohomology of all
thickenings.

% before: gal-realization-thickening-torsor-vanishing
A torsor datum also gives vanishing. When the object of the datum has the same pieces of a ladder in as many degrees as
the dimension, the cohomology of the corresponding quotient is zero.

% before: gal-realization-thickening-level-one-serre
Finiteness of the graded pieces at level $1$ is proved one property at a time for a class of modules closed under
extensions, here the finitely generated modules. For a double complex such a property passes from the cohomology of the
complexes formed by the cohomology of the columns to the total complex, through the filtration by columns.

% before: gal-realization-thickening-level-one-cech
For a category covered by finitely many cosieves, the cochains on the intersections form a double complex whose rows are
exact. Its columns give the \v{C}ech complexes of the cohomology of the intersections, so the property passes from these
\v{C}ech complexes to the cohomology of the whole category.

% before: gal-realization-thickening-level-one-graded-filtration
Over an affine open with coordinates, the graded pieces of the divided power filtration are computed through the torsor
datum of a chart. The computation uses only the multiplicative properties of the filtration, the agreement of two maps
out of the object of the datum up to one more step of the filtration, and the cochains formed from the coordinates of the
datum. So it works for any divided power functor.

% before: gal-realization-thickening-level-one-graded-forms
The coordinates of a torsor datum are additive along composites of maps. So the cochains formed from them have values in
the expected piece of the filtration and coboundaries one step deeper, and they define classes in the graded pieces.

% before: gal-realization-thickening-level-one-graded-cut
The computation goes by induction on the dimension, through the cut category, in which the first coordinate vanishes.
Along the left adjoint of the functor from the cut category, the cochains of the cut datum agree with those of the
original datum with the indices shifted by one, up to the next piece of the filtration. The cochains whose first index is
the first coordinate become coboundaries there, up to the same error.

% before: gal-realization-thickening-level-one-graded-sequence
The exact sequence of the divided power polynomial structure, transported through this adjunction, relates the graded
pieces on the category to those on the cut category. Its maps send the classes of the cochains of the original datum to
those of the cut datum, and its connecting map adds the first coordinate in front of the indices.

% before: gal-realization-thickening-level-one-graded
A torsor datum also computes the graded pieces of a filtration, when maps out of its object respect the graded pieces. In
degree $j$ the graded piece of weight $r$ is described by the forms built from the coordinates of the datum, with
coefficients of weight $r-j$ at its object.

% before: gal-realization-thickening-level-one-fibre
At level $1$ the coefficients of these forms become functions on the base change of the scheme to $W_1$. A thickening
evaluates such functions in its ring modulo the divided power ideal, and lifts of the values act on the graded pieces and
build the forms.

% before: gal-realization-thickening-level-one-forms
Since lifts are determined up to the divided power ideal, the action and the forms do not depend on their choice in the
graded pieces. The forms are cocycles, and they are compatible with sums, with the action and with restriction to
smaller opens.

% before: gal-realization-thickening-level-one-local
Over an affine open with coordinates, a chart gives a torsor datum whose ring is the ring of functions of the base
change, as its divided power ideal vanishes at level $1$. So the graded piece of weight $r$ has cohomology only in degree
$r$, and there it is free on the forms in the coordinates.

% before: gal-realization-thickening-level-one-coherent
Free modules over the opens of a cover that are related by matrices of functions on the overlaps glue to a vector bundle.
On a proper scheme over a Noetherian ring its \v{C}ech cohomology is finitely generated, by Grothendieck's finiteness
theorem.

% before: gal-realization-thickening-level-one-bundle
The forms of the charts of a coordinate cover give such data for the \v{C}ech complex of the graded piece of weight $r$
in degree $r$. Their coefficients are the matrices that express the forms of one chart through those of another on the
overlaps.

% before: gal-realization-thickening-level-one
At level $1$ the ideal $\ell W_1$ is zero, so the divided power filtration of a chart vanishes in positive degrees. The
filtration of the cohomology of thickenings then stops above the dimension, and its graded pieces are finitely generated
over $W_1(k)=k$ for a proper scheme, for every perfect field $k$ of characteristic $\ell$.

% before: gal-realization-frobenius-witt-bockstein
For a smooth projective $\mathcal O_K$-scheme, the cohomology of thickenings over the $W_\nu$ has Bockstein sequences
and finite length at the first level. By Lemma~\ref{gal-realization-frobenius-bockstein-limits} it then satisfies the
hypotheses of the flat base change for limits below.

% before: gal-realization-frobenius-flat-limit
Extending scalars to a flat algebra need not commute with inverse limits. Under the conclusions of
Lemma~\ref{gal-realization-frobenius-bockstein-limits} it does up to the power $\ell^c$, which is the source of the
factor $\ell^e$ in Lemma~\ref{gal-realization-period-ingredients}(4).

% before: gal-realization-frobenius-alpha
Combined with a base change isomorphism at each level, the limit statement gives the map $\alpha$. A compatible family
of classes in the levels comes, after multiplication by $\ell^c$, from a unique element of $A\otimes_WL$.

% before: gal-realization-period-ingredients
The existence of period and Frobenius data reduces to the following properties of the constructions
above. They are Alpöge's counterparts of the comparison theorems between crystalline, de Rham and \'etale
cohomology of Faltings, of Bhatt, Morrow and Scholze, and of Berthelot and Ogus.

% before: gal-realization-inputs
The remaining inputs are the multiplicativity and Galois equivariance of the comparison maps and the
existence of the two kinds of data. The data are assembled from the cohomologies of thickenings above,
with the bounds, Chern classes, periods and Frobenius factorization of
Lemma~\ref{gal-realization-period-ingredients}.

% before: gal-crystalline-relation-package
The crystalline vectors of a representation are its invariants in the tensor product with a ring of
crystalline periods. The comparison below is stated for every ring obtained from $A_{\mathrm{cris}}$ by
inverting $t$ and killing $t$-power torsion, with a Frobenius and a Galois action compatible with those of
$A_{\mathrm{cris}}$; Fontaine's $B_{\mathrm{cris}}$ is one.

% before: gal-crystalline-relation-package-embedding
Such a ring is the localization of $A_{\mathrm{cris}}$ at the powers of $t$. Since $t$ is invertible in
Fontaine's field $B_{\mathrm{dR}}$, it embeds there, compatibly with the Galois actions, and an identity can be
checked after this embedding.

% before: gal-crystalline-relation-package-invariants
The Galois invariants of such a ring are those of $B_{\mathrm{cris}}$. An invariant element is a quotient of an
element of $A_{\mathrm{cris}}$ by a power of $t$ that is invariant in $B_{\mathrm{cris}}$, and Fontaine's
computation applies to it.

% before: gal-crystalline-realization
The comparison with \'etale cohomology connects these groups with the Galois representation. Two of
its properties are used: a crystalline vector is determined by its classes on points over untilts, and
after a Frobenius twist it comes from a class on points of the special fibre on which $\phi^f$ acts as
pullback along the $q$-power Frobenius.

% before: gal-crystalline-splitting-cover
The finite maps $p_j$ of the correspondences do not split over $Y$, so a correspondence does not act on
points of $Y$. Over the cover constructed next every $p_j$ splits into sections, and each correspondence
becomes a sum of morphisms. The symmetric groups permuting the sections act on the cover, which allows
averaging back to $Y$.

% before: gal-crystalline-norm-forms
Over the special fibre $p_j$ need not be \'etale, so a point of $Y$ there has no well-defined multiset of
Hecke translates. The norm form of a function along $p_j$ is a polynomial over the base, and wherever
$p_j$ splits it gives the multiset of values of the function at the points of the fibre.

% before: gal-crystalline-norm-identity
The Eichler--Shimura relation is assumed only at the generic points of the special fibre. In terms of
norm forms it becomes an identity of polynomials over an affine chart modulo each minimal prime over
$\varpi$. Since the special fibre is reduced, the identity then holds after every base change that kills
$\varpi$.

% before: gal-crystalline-matching-products
Over a field, equal products of the polynomials $\Phi_c$ force equal multisets of vectors, since $\Phi_c$
determines $c$ and the polynomial ring has unique factorization. Over a ring with zero divisors the
matching of two multisets is built one pair at a time, and each step turns a congruence of products
modulo a power of $w$ into a congruence of vectors modulo a smaller power.

% before: gal-crystalline-matching-forms
The matching argument uses four properties of products of the forms $\Phi_c$, which hold over every ring
because each $\Phi_c$ is monic of degree one in $X$. Products are compatible with base change, a common
factor cancels, the degree in $X$ counts the factors, and replacing $X$ by the linear form of one member
kills the product.

% before: gal-crystalline-matching-gauss
Over a ring with zero divisors a product of polynomials can vanish while no factor does. A constant times
a product of polynomials, each with a coefficient equal to $1$, vanishes only if the constant does. This
form of Gauss's lemma is applied to the linear forms $\sum_k(x_k-y_k)T_k$, after a common factor has been
taken out of each.

% before: gal-crystalline-matching-factor
To take a common factor out of the linear form $\sum_k(x_k-y_k)T_k$ and leave a coefficient equal to
$1$, one coordinate difference must divide all the others. This holds when each coordinate that divides
a given power of $w$ is comparable for divisibility with every other coordinate, which the localization
to idempotents arranges.

% before: gal-crystalline-matching-step
The inductive step matches one member $x$ of the first multiset with a member of the second, modulo
$w^r$. If no member matched, the product of the common factors of the forms $\sum_k(x_k-y_k)T_k$ would
divide $w^{rN}$ and, by Gauss's lemma, be divisible by $w^{r(N+1)}$. For this reason the congruence of
products is assumed modulo $w^{r(N+1)}$.

% before: gal-crystalline-matching-induction
Removing a matched pair and cancelling its forms leaves a congruence for the smaller multisets modulo
$w^r$. Starting from $w^{(N+1)^N}$ and dividing the exponent by $N+1$ at each of at most $N$ steps, the
induction matches the two multisets modulo $w$. In a valuation ring the divisibility hypotheses hold
automatically, since any two elements are comparable for divisibility.

% before: gal-crystalline-matching-piece
In a test algebra two elements are comparable for divisibility only after a decomposition by idempotents.
Over the factor $B/(1-e)B$ of $B$ cut out by an idempotent $e$, divisibilities multiplied by $e$ become
divisibilities, and the induction applies there.

% before: gal-crystalline-matching-local
The divisibility property of a test algebra compares two elements on the two pieces of a decomposition by
an idempotent. Near a prime ideal, finitely many comparisons can be made at once, by multiplying
idempotents outside the prime. The comparisons needed are those of the differences of coordinates with
each other and with the powers of $w$.

% before: gal-crystalline-matching-cover
The idempotents found near different maximal ideals may match the two families by different
permutations. Finitely many of them generate the unit ideal, and ordering them gives a decomposition of
$1$ into orthogonal idempotents indexed by the permutations.

% before: gal-crystalline-matching
Equal norm forms over a ring with nilpotent elements need not come from equal multisets of values. Over
a test algebra modulo $\ell^\flat$ they do after a large power of Frobenius and a decomposition, because
the test algebra is perfect and has the divisibility property of a valuation ring.

% before: gal-crystalline-points-relation
The norm identity and the matching lemma together turn the Eichler--Shimura relation into an identity
between points on test algebras. Reduced to the special fibre and twisted by a power of Frobenius, the
weighted Hecke translates of a point of the splitting cover cancel locally.

% before: gal-crystalline-relation
This is the counterpart of Lemma~\ref{gal-eichler-shimura-cohomology} at a place over $\ell$, with the
crystalline Frobenius in place of the geometric Frobenius. A crystalline vector is moved to a class on
points of the special fibre, where the identity between points applies, and moved back by the
comparison of Theorem~\ref{gal-crystalline-realization}, in any crystalline period package, Fontaine's
$B_{\mathrm{cris}}$ among them. The extra power $\phi^{fN}$ comes from the Frobenius twist in that comparison
and from the matching lemma.

% before: gal-tower-cohomology
The properties of a seed come from four properties of the cohomology of a smooth projective tower.
They are the existence of a semisimple Galois piece, Frobenius eigenvalues from the congruence
relation, crystallinity and Hodge--Tate weights at places of good reduction over $\ell$, and the
crystalline Frobenius.

% before: gal-pack-seed
A pack gives a seed by taking a piece of the cohomology of its tower and twisting it by the $\ell$-adic
character of the Hecke character. The twist moves the Frobenius eigenvalues and the Hodge--Tate weights
to those predicted by $U$. Lacunary weights at one place over $\ell$ give the last condition of a
seed.

% before: gal-pel-datum
The Hecke towers of the packs are Shimura varieties of PEL type for unitary groups of signature
$(1,n-1)$ at one infinite place and definite elsewhere. They are proper because the unitary group is
definite somewhere, so the smooth projective statements above apply to them.

% before: gal-holomorphic-lefschetz
To show that a system of Hecke eigenvalues occurs in a given Hodge degree, one computes the alternating
trace of a Hecke operator on coherent cohomology. For a correspondence with isolated nondegenerate fixed
points, the holomorphic Lefschetz formula gives this trace as a sum over the fixed points.

% before: gal-plan
A plan records the group-theoretic data from which a pack is built: a PEL datum, a level, Hecke
operators with their eigenvalues, a test function that detects these eigenvalues spectrally, and the
congruences that the Eichler--Shimura relation will turn into Frobenius eigenvalues. Separating the plan
from the pack keeps the automorphic input apart from the geometry.

% before: gal-powers-situation
The congruence relation of Proposition~\ref{gal-ks-towers}(3) is proved on an integral model at a
place $w$ over a prime at which the PEL datum is also of simple type. The model is the moduli space of
the datum over a discrete valuation ring above $w$, with an extra level structure at the other places
over the same prime of the imaginary quadratic field. That level structure is \'etale: by the Kottwitz
condition, the part of the Lie algebra over this prime lies at $w$, so the torsion at the other places
over it is \'etale.

% before: gal-powers-hecke
The Hecke operators at $w$ that enter the congruence relation come from isogenies whose kernel meets the
torsion over $\mathfrak u$ in a $j$-dimensional subspace of the $w$-torsion. Away from $p$ these
isogenies are classified by a finite \'etale scheme over the model, and the correspondence of type $j$
extends this scheme over the whole model.

% before: gal-powers-hecke-exist
The correspondences of type $j$ exist on the integral model, as the normalization of the classifying
scheme over the locus where $p$ is invertible. They pass to the model with extra level structure,
because an isogeny of type $j$ is injective on the torsion at the other places over $\mathfrak u$, and
to any scheme over the model by base change, with a fibre count given by subspaces.

% before: gal-powers-subspace-identity
At an ordinary point of the special fibre, the reduction of an isogeny of type $j$ depends on a
$j$-dimensional subspace of an $n$-dimensional space only through its image modulo a line, and a
subspace containing the line contributes a Frobenius twist. The Eichler--Shimura relation then comes
down to a count of subspaces, which telescopes.

% before: gal-powers-generic-ordinary
The count of the previous lemma applies at points where the $w$-torsion has as many points in
characteristic $p$ as the dimension of the moduli space allows. The generic points of the special fibre
are such points, by the stratification of Harris and Taylor.

% before: gal-powers-special-isogenies
Over an algebraically closed field of characteristic $p$, an isogeny of type $j$ is determined, with its
target, by its kernel over $\mathfrak u$. Two such isogenies whose kernels have the same points have
the same target when the kernels are \'etale. When one kernel has an extra infinitesimal part of order
$q$, the targets differ by the $q$-power Frobenius.

% before: gal-powers-reading
At a point with values in a valuation ring whose reduction is $w$-ordinary, the $w$-torsion is an
$n$-dimensional space and its points reducing to zero form a line. The isogenies of type $j$ at the
point correspond to the $j$-dimensional subspaces, and the previous lemma reads off their reductions
from the subspaces modulo the line.

% before: gal-powers-eichler-shimura
On the fibre powers of the universal abelian scheme, the Hecke correspondences act through the isogeny on
each factor and on the level structure. A point of the fibre power is determined by these data, so the
reading at $w$-ordinary points passes to the fibre powers, and the count of
Lemma~\ref{gal-powers-subspace-identity} gives the Eichler--Shimura relation at the generic points of
the special fibre.

% before: gal-count-situation
Item (4) of Proposition~\ref{gal-ks-towers} detects Hodge degrees by counting fixed points of Hecke
correspondences on the complex points of the Shimura variety and comparing the count with a trace on a
definite unitary group. The count uses only the group of the PEL datum: a central simple algebra with
an involution of the second kind, of signature $(1,n-1)$ at one infinite place and definite at the
others, together with its local data at $p$.

% before: gal-count-forms
The complex points of the Shimura variety are described by forms, Hermitian elements of the algebra
that are locally trivial, together with positive lines at the indefinite place. A Hecke correspondence
fixes such a point when an element of the group of the form fixes its line and moves its adelic
coordinate into the support of the test function, and the complex fixed-point sum weights these
elements by a function of their minimal polynomial and eigenvalue.

% before: gal-count-pairs
Grouping fixed points by the conjugacy class of the fixing element turns the fixed-point sum into a sum
over pairs, each weighted by an orbital sum of the test functions. The comparison is with the groups of
definite forms, which are compact at infinity, and the traces of Hecke operators on functions on these
groups are what the spectral condition of a plan controls.

% before: gal-count-positive-line
A regular element of the group of a form has $n$ distinct eigenvalues at the indefinite place, where
the Hermitian form has a single positive direction. So it fixes at most one positive line, and the
eigenvalue at a fixed point is well defined.

% before: gal-complex-majorant
The fixed-point sum is a sum over infinite sets, and the next lemmas show that only finitely many of its
terms are nonzero. At the indefinite place the Hermitian form of a form is not definite, so it does not
bound the elements of its group. An element that fixes a positive vector preserves a positive definite
form built from that vector, and this restores the bound.

% before: gal-complex-archimedean
At the infinite places other than $v$ the Hermitian form of $\beta$ is definite, and at $v$ the previous
lemma supplies a definite form preserved by the elements that fix $z$. The entries of these elements are
then bounded in terms of their similitude factor.

% before: gal-complex-lattice
Bounded denominators put the relevant elements of $C$ in a finitely generated group, and the bound at the
infinite places then leaves finitely many of them. This works because the complex representations at the
infinite places together embed $\mathbf R\otimes C$.

% before: gal-complex-fixed-finite
The denominators of an element of the group of a form are bounded at the finite places by the supports of
$f$ and $f_p$, and its size at the infinite places by the positive vector it fixes. Together these leave
finitely many elements fixing a given positive vector with a nonzero summand.

% before: gal-complex-separable
The finiteness of the points needs the elements with a nonzero summand to have separable minimal
polynomials of degree $n$. Regularity at one prime makes their centralizers commutative, and the definite
form preserved at $v$ rules out nilpotent parts.

% before: gal-complex-separable-degree
The proof above obtains the degree from an eigenspace argument over an algebraic closure. Knowing the
centralizers of diagonalizable matrices gives a shorter route.

% before: gal-complex-polynomials
The first reduction is to finitely many minimal polynomials. Their coefficients have bounded
denominators, now uniformly in the point, because the characteristic polynomial does not change under
conjugation, and their roots at the infinite places all have absolute value the square root of the
similitude factor.

% before: gal-complex-cosets
For a regular semisimple element, the elements conjugating it into a compact set lie in finitely many
cosets modulo its centralizer. At almost all primes this is Kottwitz's lemma, and at the remaining primes
and at $p$ it is the properness of the orbit map.

% before: gal-complex-torus
The centralizer of a regular semisimple element of the group of a form is a commutative algebraic group,
whose adelic points have finitely many classes modulo its rational points and an open subgroup. This turns the cosets of the
previous lemma into finitely many points.

% before: gal-complex-points-conjugate
For a single conjugacy class of the group of a form, the two previous lemmas and the uniqueness of the
positive line leave finitely many points. The fixed positive vector is moved to a standard one by the
conjugating element and the centralizer.

% before: gal-complex-classes-local
Two elements with the same minimal polynomial are conjugate in $C^\times$, and the conjugating element
gives a unit of the centralizer fixed by the involution, which decides whether they are conjugate in the
group of the form. At almost all primes this unit is a local norm up to a scalar, by Kottwitz's lemma and
local class field theory.

% before: gal-complex-classes-units
The local norm conditions leave finitely many global classes. After the scalar is made rational, this is
a finiteness for the fixed elements of a number field with an involution modulo norms, which follows from
the finiteness of the class group and of the units modulo squares.

% before: gal-complex-classes
With the two previous lemmas, a minimal polynomial accounts for finitely many conjugacy classes of the
group of a form. The class of an element is determined by the class of its unit and the sign of a
rational number.

% before: gal-complex-points-finite
The three reductions combine: finitely many minimal polynomials, finitely many conjugacy classes for each,
and finitely many points for each class. This is the second finiteness used in the proof of the next
lemma.

% before: gal-count-complex
Since the positive line is unique, each regular conjugacy class contributes once, with its eigenvalue
on that line, and the points it fixes assemble into an orbital sum. The fixed-point sum becomes a sum
over the classes of pairs with a marking.

% before: gal-ratio-integral
Two pairs with the same $\gamma$ and $c$ are compared through the ratio of their forms, prime by prime.
Away from finitely many primes the comparison takes place in the integral structure that a basis of $C$
defines at each prime, and it uses the involution $x\mapsto\beta x^*\beta^{-1}$ attached to a form
$\beta$.

% before: gal-ratio-centralizer
The ratio of the forms of two pairs with the same $\gamma$ commutes with $\gamma$. When the minimal
polynomial of $\gamma$ is separable of degree $n$, the centralizer of $\gamma$ is a product of number
fields, the involution of a form restricts to an automorphism of it, and both are compatible with
completion at a prime.

% before: gal-ratio-order
At all but finitely many primes the lattice spanned by the basis is a ring stable under the involution,
it contains a given form and its inverse, and its elements are integral over $\mathbf Z_\ell$. Each of
these conditions concerns finitely many rational numbers.

% before: gal-ratio-split
Conjugates of $\gamma$ that are integral at a prime are compared through a matrix description of the
completion of $C$. At almost all primes the integral structure becomes the matrices over the integers of
the completion of $F$, because an order of $C$ is maximal and $C$ is split at almost all places.

% before: gal-ratio-semilocal
The integers of the completion of $F$ at a prime have finitely many maximal ideals. Over such a ring two
integral matrices with the same separable characteristic polynomial are conjugate by an integral
invertible matrix, since each of them has a cyclic vector modulo every maximal ideal.

% before: gal-ratio-lattice
At a prime, the element $k$ that carries the form of a pair to a scalar conjugates $\gamma$ into the
integral structure. Conjugacy of integral matrices splits $k^{-1}$ into a factor commuting with $\gamma$
and an integral unit, so only an integral unit has to be controlled.

% before: gal-ratio-unit-norm-finite
An integral unit fixed by the involution of a form has to be written as a norm for that involution. The
argument works for any ring with an involution for which some $\omega-c(\omega)$ is a unit, first for
finite rings. There a version of Hilbert's Theorem 90 and a count of the units show that the fixed units
are the norms.

% before: gal-ratio-unit-norm-hensel
Norms then lift from a finite quotient to a Henselian pair. Units congruent to $1$ are norms by Hensel's
lemma for a quadratic equation, and a general fixed unit is a norm modulo the ideal.

% before: gal-ratio-unit-norm
Over the integers of a local field of odd residue characteristic the norm form is $a^2-db^2$ with $d$ a
unit. Adjoining a square root of $d$ produces a ring with an involution to which the previous lemma
applies.

% before: gal-ratio-unit-norm-count
The proof above passes through the general statement for Henselian pairs. For the form $a^2-db^2$ one
can instead count squares in the residue fields.

% before: gal-ratio-local-norm
The involution of a form changes the sign of an element $j$ of $F$ with $j^*=-j$, so the centralizer of
$\gamma$ is generated over its fixed part by $j$. A fixed unit is then the norm of $a+bj$ as soon as it
equals $a^2-j^2b^2$, which the previous lemma provides at almost all primes.

% before: gal-ratio-local
These local results give the constraint on the ratio of the forms of two pairs with the same $\gamma$
and $c$. Outside a finite set of primes that depends only on the test function and on one of the pairs,
a nonzero term of the orbital sum of the other pair makes the ratio a scalar times a norm for the
involution of the first form.

% before: gal-count-finite-integral-degree
The bound on denominators below holds for orbital sums whose factor at $p$ is a twisted orbital integral
over the unramified extension of $\mathbf Q_p$ of degree $r$. In degree $r$ the element at $p$ is an element
whose norm is conjugate to the image of $\iota_p(\gamma,c)$, and the test function at $p$ is invariant under
conjugation twisted by the Frobenius. In degree $1$ these are the orbital sums already defined.

% before: gal-count-finite-integral
The first step towards finitely many classes bounds the denominators of $\gamma$. Away from $p$ the
support of $f$ bounds those of its characteristic polynomial, and at $p$ the support of $\varphi$ bounds the
entries of the norm, and so the eigenvalues of $\gamma$. One integer then clears the denominators of
$\gamma$ and $\gamma^{-1}$ for every pair with a nonzero orbital sum.

% before: gal-count-finite
Rearranging the sums over classes of pairs needs that only finitely many classes contribute. The
support of the test functions bounds the minimal polynomial of a pair with a marking, and for a fixed
polynomial the forms differ by elements of a commutative algebra that are constrained locally.

% before: gal-count-definite-signs
Passing from a pair with a marking to a definite pair multiplies the form by a Hermitian element of the
field generated by the pair. This element must change the sign at the eigenlines other than the marked
one and keep the form locally trivial at the finite places, which Hilbert reciprocity allows when the
number of sign changes is even.

% before: gal-count-definite-move
For odd $n$ the number $n-1$ of signs to change is even, so one element of the field generated by
$\gamma$ moves every pair with a chosen marking to a definite pair. The move does not change the form at
the finite places away from $p$, so the minimal polynomial and the orbital sums are kept.

% before: gal-count-weighted
Summing the bijection over the $n$ possible marked roots turns the complex fixed-point sum into a sum
over definite pairs, weighted by the sum of the function over all roots at $v$. The pairs with a
marking and a reducible minimal polynomial have no counterpart, and their contribution is assumed to
vanish.

% before: gal-definite-form-similitude
The trace of a Hecke operator on a definite form is a finite sum because the group of the form has
finitely many double cosets in $G(\mathbf A_f^p)\times G_p$, each with a finite stabilizer. Both facts are
proved through the similitude group of the form. It contains the group of the form as its elements with
positive similitude factor, and its points away from $p$ and at $p$ can be compared with those of $G$.

% before: gal-definite-form-class-number
The finiteness of the double cosets comes from Borel's finiteness of class numbers for the similitude
group of a Hermitian element, and that of the centralizers in Lemma~\ref{gal-complex-torus} from the same
theorem for the elements of this group that commute with a given element. It covers the points away from $p$
and at $p$ together by finitely many translates of an open subgroup by rational points.

% before: gal-definite-form-padic
The group $G_p$ records only the components at the places $w_i$ over $u$. Since $p$ splits in $K_0$,
the components of $\mathbf Q_p\otimes C$ at the conjugate places are recovered through the involution, so
the $p$-adic points of every similitude group are described by $G_p$.

% before: gal-definite-form-classes
For a definite form, the maps $\iota$ and $\iota_p$ extend to the points of its similitude group away
from $p$ and at $p$, by conjugation by $x$ and by the homomorphisms $\rho_i$. Borel's theorem then covers
$G(\mathbf A_f^p)\times G_p$ by finitely many double cosets. A level $K'$ commensurable with $K_p$ adds
finitely many coset representatives.

% before: gal-definite-form-denominators
The finiteness of the stabilizers rests on bounds for the coordinates of an element of the group of a form
in a basis of $C$. Where a test function is nonzero at the image of the element, the coordinates have
bounded denominators: away from $p$ because the support of a test function is compact, and at $p$ by the
description of $\mathbf Q_p\otimes C$ through the homomorphisms $\rho_i$.

% before: gal-definite-form-infinite
At the infinite places the Hermitian matrices of a definite form are definite, so an element of its group
stretches the length of every vector by at most a constant times the square root of its similitude
factor. This bounds the coordinates of the element by the same quantity.

% before: gal-definite-form-stabilizers
In a stabilizer, the denominator bound also bounds the similitude factor, since a fixed integer times its
inverse is a positive integer. The two bounds then confine the coordinates of the elements to a finite
set of rational numbers.

% before: gal-count-definite-trace
On a definite group the automorphic forms of a given weight are functions on a finite set, and the trace
of a Hecke operator is a sum over conjugacy classes of the character of the weight times an orbital sum.
When the weight of the definite pairs is this character, their weighted sum is the definite trace.

% before: gal-count-weights
At a fixed point the correspondence acts on the tangent spaces of the ball and of the fibres of the
abelian scheme, and its Lefschetz term depends only on the eigenvalues there. These eigenvalues are
given by the roots of the minimal polynomial and by the multiplications $\alpha$ on the fibres, which is
how the Lefschetz weights are defined.

% before: gal-count-root-weights
A combination of the Lefschetz weights for several multiplications $\alpha$ on the fibres isolates the
character of the representation of weight $\lambda$, placed in the Hodge degrees given by its
exponents. This is an identity between polynomials in the roots and in $y$.

% before: gal-count-group-side
The weighted count, the identity between root sums and the trace on definite groups combine into the
group side of the Lefschetz formula. A combination of complex fixed-point sums for the Lefschetz weights
equals the definite trace times the generating function of the Hodge exponents.

% before: gal-fixed-situation
The fixed points of the Hecke correspondences are found through the complex uniformization of the
Shimura variety, which writes each abelian variety of the moduli problem as a complex vector space
modulo the image of a lattice in the rational space of the PEL datum. The action of $\mathcal O_E$ on
that complex space is fixed once, with the multiplicities that the Kottwitz condition prescribes on the
Lie algebra.

% before: gal-fixed-uniformization
A complex point of the Shimura variety is the quotient of $\mathbf C^q$ by the image of a lattice of $V$
under a period. The lattice records the finite adelic coordinates of a point of a form, and the period
records its positive vector. Isomorphisms between such points, the Hecke correspondences and the local
coordinates on the Shimura variety are all read off from this description.

% before: gal-fixed-stabilizer
An element of the stabilizer of a point of level $N$ acts on the abelian variety of the corresponding
point as an automorphism that preserves the polarization and fixes the $N$-torsion. For $N\geq3$ such an
automorphism is trivial, so the sums over points below carry no stabilizer weights.

% before: gal-fixed-base
A fixed point on the fibre powers lies over a point of level $N'$ whose images under the change of level
and the Hecke correspondence agree. Through the uniformization, this equality is an isomorphism between
two lattices of $V$, and such an isomorphism is an element of the group of the form.

% before: gal-fixed-torus
Above a fixed point of the Shimura variety the fibre is a power of a complex torus, on which the
correspondence acts through a linear map that preserves the lattice. The fixed points of such a map are
counted by a determinant.

% before: gal-fixed-eigenvalues
The contribution of a fixed point is computed from the eigenvalues of the linear maps that the
correspondence induces on the tangent space of the Shimura variety and on the fibre. These eigenvalues
are the roots of the minimal polynomial of the element of the group, distributed over the embeddings of
$E$ by the multiplicities of the Kottwitz condition, and they assemble into the function $L_{\alpha,t}$.

% before: gal-fixed-fibre
Above one fixed point of the Shimura variety, the fixed points on the fibre powers are found in the
local coordinates of the uniformization, in which the correspondence is linear. Their Lefschetz terms
add up to the Lefschetz weight at the minimal polynomial and the eigenvalue of the element of the group.

% before: gal-fixed-level
The fixed points give a sum over points of level $N'$ weighted by the indicator of the coset
$h^{-1}K_N$, while the group side uses points of level $N$ and the indicator of the double coset
$K_Nh^{-1}K_N$. When the stabilizers are trivial, the two sums differ by the number of cosets of
$K_{N'}$ in $K_N\cap h^{-1}K_Nh$.

% before: gal-fixed-points-isolated
The Lefschetz formula needs finitely many fixed points. A fixed point at which the correspondence has
no eigenvalue $1$ on the tangent space is isolated, by Nakayama's lemma applied to the functions that
cut out the fixed points. On a scheme of finite type over $\mathbf C$ isolated fixed points are finite in
number.

% before: gal-fixed-points
The geometric side is the Lefschetz formula for a Hecke correspondence composed with multiplication by
$\alpha$ on the fibre powers. Its fixed points lie over the points of the Shimura variety that are fixed
by elements of the groups of forms, and the Lefschetz terms above each of them add up to the Lefschetz
weight.

% before: gal-ks-towers
The geometric input about the PEL Shimura varieties is collected here. The towers live on fibre powers
of the universal abelian scheme, whose cohomology realizes the automorphic forms of weight $\lambda$.
Items (1)--(3) are the properties of the integral models used in
Proposition~\ref{gal-tower-cohomology}, and item (4) turns the spectral condition of a plan into
occurrences in Hodge cohomology through the Lefschetz formula.

% before: gal-plan-pack
With the towers of Proposition~\ref{gal-ks-towers}, a plan becomes a pack. The Hodge data and the
congruences pass from the plan to the pack place by place.

% before: gal-identification-isometry
The plans compare the group of a PEL datum, whose Hermitian form is not the standard one, with the
definite unitary group of $U$. In odd rank the two forms become isometric over the finite adeles, through
an explicit block matrix built from a representation of $\gamma^{-1}$ by a sum of two norms. Integrality
at the good primes keeps this isometry compatible with the levels there.

% before: gal-identification
The isometry gives a map from the similitude group of the standard space onto the finite adelic points of
$G$ away from $p$. Levels, Hecke operators and test functions on the definite side are transported along
this map, and item (5) identifies the rational points of the definite group with the same finite adelic
points as $G$.

% before: gal-trace-vector
The automorphic input on the definite side is a vector-valued form on the similitude group. It transforms
under a supercuspidal type at $w_0$, is spherical with the Satake parameters of $U$ at the split places
outside a finite set, and has a central character given by an algebraic Hecke character; its weight is
the weight that the plan carries.

% before: gal-trace-seen-operators
The Hecke operators of a plan are correspondences of level $N$, given by a double coset and a pair of
levels. Two kinds are used: deck transformations of the level structure, and the operators $T_{w,i}$ at
places free for the vector, whose values on $U$ are the normalized elementary symmetric functions of its
Satake parameters.

% before: gal-trace-auxiliary-prime
The test function needs auxiliary places at which its support consists of regular elements. They are
found over a large prime that splits in $\mathbf Q(\omega)$, away from a given finite set and from the
places of the vector, and the places of $E$ over it come in conjugate pairs.

% before: gal-test-function
The test function is built from the level of the vector, the supercuspidal type at $w_0$ transported by
the inverse transpose, and Iwahori subgroups at the auxiliary places. Convolution with it projects the
forms of the definite group onto those that transform under this type.

% before: gal-test-function-type
The type inherits its properties from the vector and from the identification. Item (4) makes the test
function cuspidal at $w_0$, and item (5) separates its behaviour at $w_0$ from that at the other places.

% before: gal-iwahori-operator
At the auxiliary places the test function needs an inverse, in the Hecke algebra of the Iwahori subgroup,
of the characteristic function of a double coset. This inverse comes from the action of double cosets on
the vectors fixed by the Iwahori subgroup in an arbitrary representation. The operator of a double coset
adds up the translates of a fixed vector by representatives of its cosets. Alpöge works with the subgroups
of the Iwahori subgroup cut out by a condition on the residues of the diagonal entries, from the pro-$p$
Iwahori subgroup up to the Iwahori subgroup itself, and the lemmas below hold for all of them.

% before: gal-iwahori-operator-product
Two such operators compose to the operator of the product double coset when the products of
representatives run once through its cosets. Containment and a count of cosets are enough to check this.

% before: gal-iwahori-operator-count
The count of cosets in a product double coset reduces to two conditions on how the two elements
conjugate elements of the subgroup. For $\operatorname{GL}_n$ both conditions come from the Iwahori
factorization.

% before: gal-iwahori-operator-normalizer
An element that normalizes the subgroup has a double coset consisting of a single coset. Its operator is
the translation by the element, which is invertible and can be split off from any other operator.

% before: gal-iwahori-operator-factorization
Every element of the Iwahori subgroup factors into an upper unitriangular, a diagonal and a lower
unitriangular matrix. Conjugation by a dominant diagonal matrix keeps the upper factor integral, and
conjugation by its inverse keeps the lower factor in the Iwahori subgroup.

% before: gal-iwahori-operator-dominant
For two dominant diagonal matrices the conditions of Lemma~\ref{gal-iwahori-operator-count} follow from
the factorization. So the operators of their double cosets multiply, and it suffices to invert the
operators of a few dominant matrices from which the others are built.

% before: gal-iwahori-operator-simple
These generators are reached through permutation matrices. For the transposition of two adjacent
indices, the double coset has one coset for each residue class, and the square of its operator is a
multiple of the identity plus terms that the operator divides on either side. This makes the operator
invertible when the residue characteristic is invertible in the coefficients, and for the Iwahori
subgroup it is the usual quadratic relation.

% before: gal-iwahori-operator-permutation
A permutation is a product of adjacent transpositions, one for each of its inversions. Each step
multiplies the number of cosets by the size of the residue field, so the operators multiply along such a
product, and the operator of every permutation matrix is invertible.

% before: gal-iwahori-operator-fundamental
The matrix with $m$ diagonal entries equal to the uniformizer and the others equal to $1$ does not
normalize the Iwahori subgroup. Composed with a cyclic permutation matrix it does, and this reduces its
operator to that of a permutation matrix.

% before: gal-iwahori-operator-invertible
Every dominant diagonal matrix is a unit diagonal matrix times a scalar times a product of the matrices of
Lemma~\ref{gal-iwahori-operator-fundamental}. The multiplicativity of
Lemma~\ref{gal-iwahori-operator-dominant} then makes the operator of its double coset invertible. The factorization of
the Iwahori subgroups also shows that these operators commute with the passage between two of the groups $I_C$.

% before: gal-test-function-iwahori
At the auxiliary places the test function is modified by an Iwahori double coset whose powers consist of
regular semisimple elements. The three properties below are local facts about
$\operatorname{GL}_n$ over a non-archimedean local field.

% before: gal-test-function-powers
Translating the double cosets of Lemma~\ref{gal-test-function-iwahori} to the auxiliary places, and
multiplying by a central scalar to make them integral, gives a function $u$ that is invertible on the image
of $e$. Its powers applied to $e$ are supported on elements whose components at the auxiliary places have
separable characteristic polynomial.

% before: gal-hecke-type-functions
The relations of Lemma~\ref{gal-test-function-hecke} are proved by letting functions on $G(\mathbf A_f^p)$
act on functions that transform under $K$ by the character $\psi$. The function $\phi$ coming from $U$
transforms this way, and so do the translates of $e$, through which an action on such functions becomes a
convolution with $e$.

% before: gal-hecke-double-cosets
A double coset of a subgroup generated by an element at one split place can be summed over at that place,
provided the subgroup contains the local subgroup onto which it projects. The argument uses only a
projection to the place and a section of it whose image commutes with the kernel. This reduction is used
for the group $K$ of the type and for the group of a Hecke operator of kind (2)(b).

% before: gal-hecke-test-localization
A projection with such a section is a localization of the group. The action of test functions at a split
place is computed for any localization whose section is invisible to the projection to $w_0$; the
localizations at two split places over different places of $F$ have this property.

% before: gal-hecke-test-action
The test functions $e_h$ with $\operatorname{pr}_{w_0}(h)=1$ transform on both sides by $\psi^{-1}$, and
convolution with $e$ on the right fixes them. On functions of type, $e_d$ for a central $d$ is right
translation by $d$, and the test function of an element in the image of a localization apart from the one
at $w_0$, times $d$, is a double coset operator of the localizing group followed by this translation.

% before: gal-hecke-convolution
An identity between the actions of two functions on all functions of type gives the same identity between
their convolutions with $e$, because the translates $x\mapsto e(x^{-1}g)$ are of type. Item (2) of
Lemma~\ref{gal-test-function-hecke} is obtained this way from the actions computed below.

% before: gal-hecke-seen-shape
The group of a seen Hecke operator lies in $K$, where $\psi$ is trivial. The element of an operator of kind
(2)(b) is a matrix at a split place prime to $p$, away from $w_0$ and the auxiliary places, times a central
element, and the group of the operator contains the integral matrices at that place. So each operator can
be compared with a test function.

% before: gal-hecke-seen-action
On functions of type, a seen Hecke operator of kind (2)(a) acts as the identity, and one of kind (2)(b) acts
as its index times the test function of the inverse of its element. In the second case the cosets of the
operator are grouped into double cosets, each counted as often as the index, and the double coset sum is
then computed at the split place of the operator.

% before: gal-hecke-satake-eigenvalue
At the split place of an operator of kind (2)(b), the double coset sum applied to the function of the vector
after the automorphism $\theta$ is given by the Satake parameters of $U$ and by the value of $\kappa$ at that
place. The inverse transpose in $\theta$ turns the double coset of $t_i^{-1}$ into that of $t_i$, which is
the $i$-th Hecke operator.

% before: gal-test-function-hecke
The Hecke operators of the plan preserve the image of the test function, commute there, and act on the
form coming from $U$ by the prescribed values. These are the relations needed to reduce the traces to
linear algebra.

% before: gal-trace-eigenform
The form of the vector, transported to the definite group, is a common eigenvector of the Hecke operators
fixed by the test function. Its representation at infinity is the inverse transpose of that of the vector,
and its weight is the weight $\lambda_V$ of the plan.

% before: gal-trace-linear-algebra
The nonvanishing of the trace reduces to linear algebra on the finite-dimensional spaces of forms of the
definite group. A common eigenvector in the image of $e$ survives in a generalized eigenspace, and an
invertible operator on a nonzero space has a power with nonzero trace.

% before: gal-trace-cuspidal-vanishing
The fixed-point count of Proposition~\ref{gal-ks-towers} contains terms at elements that are not regular
elliptic. For functions that are cuspidal at one place and have regular support at another these terms
vanish, by the vanishing of orbital integrals of cusp forms at non-elliptic elements.

% before: gal-trace-nonzero
With these ingredients the spectral condition of a plan can be produced. The auxiliary prime is chosen
first, the test function is a power of the regular operator $u$ applied to the type idempotent, and the
trace is nonzero because the form transported from $U$ contributes with the eigenvalues $\vartheta$.

% before: gal-plans-exist
It remains to produce plans from $U$. The definite unitary group of $U$ and the PEL group have
isomorphic finite adelic points away from the auxiliary prime, so Hecke operators and test functions can
be transported between them, and the supercuspidal type at $w_0$ makes the trace of the test function
nonzero.

% before: gal-packs-exist
Plans and the passage from plans to packs together give the packs needed for the seeds. The three items
match the three conclusions of Theorem~\ref{gal-seed-representations}.

% before: gal-weakly-attached
The construction of the Galois representations proceeds by extending the class of cuspidal
representations for which they are known. The basic notion is a continuous semisimple representation
with the right Frobenius eigenvalues at almost all places. Everything else is a condition on such a
representation at further places.

% before: gal-spherical-conditions
At a place where the automorphic representation is spherical, four conditions describe the expected
local behaviour. They are unramifiedness with the right Frobenius eigenvalues away from $\ell$, and
de Rham, crystalline and Frobenius conditions over $\ell$. The construction keeps track of all four.

% before: gal-scope
The construction starts with cuspidal representations over CM fields with enough splitting and with
conjugate self-dual polarization, and extends from there. The three scopes name the stages. Existence
and the four conditions are statements about each scope.

% before: gal-plain
Local-global compatibility at a place where the local component is ramified is reduced to the
spherical case by base change. This works when the local component is induced from characters in a
way that base change can make unramified. Plain representations are those.

% before: gal-generic-monodromy-free
The compatibility clauses of Definition~\ref{gal-compatible} are stated in terms of the Weil--Deligne
representation $\operatorname{rc}(\sigma^\flat)$, while the propositions that prove them are stated in
terms of the supercuspidal support of $\sigma$. For a generic $\sigma$, the laws translate one into the
other.

% before: gal-solvable-disjoint
The proof of Proposition~\ref{gal-plain-compatibility} carries $\pi$ along a tower of CM fields with
solvable Galois group, one cyclic step of prime degree at a time. Cyclic base change can fail to be
cuspidal, and the tower avoids this by staying disjoint from an open subgroup on which the trace of the
Galois representation does not vanish.

% before: gal-solvable-cover
Disjointness from an open subgroup can be forced by splitting conditions at finitely many places. By
the Chebotarev density theorem, the decomposition groups at finitely many places meet every coset of an
open subgroup.

% before: gal-solvable-abelian-real
The local behaviour of the tower at the place under $A$ is prescribed through an abelian extension of a
totally real field. Its existence is the Grunwald--Wang theorem on global characters with given local
components, applied to the characters of a finite abelian quotient of the decomposition group.

% before: gal-solvable-compositum
Adjoining a totally real solvable extension to a CM field gives a CM field. If the extension splits at
the places of Lemma~\ref{gal-solvable-cover}, the result is disjoint from the given open subgroup, so
the totally real fields of Lemma~\ref{gal-solvable-abelian-real} yield CM fields along which base change
can be taken.

% before: gal-solvable-cm-extensions
The lift uses two CM extensions. The first keeps the completion at the place under $A$ and makes that
place split over the maximal totally real subfield. The second needs this splitting, and cuts out a
prescribed open subgroup of the decomposition group.

% before: gal-solvable-character-kernel
The subgroup cut out by the second extension must contain a given element of nonzero degree, while the
characters of the supercuspidal support become unramified over the extension. Only its elements of
inertia are constrained, so a large unramified part keeps it open.

% before: gal-solvable-base-change
Along a tower, the lifted representation keeps the local properties of cyclic base change, and it also
carries a supercuspidal support made of characters to their composites with the norm. The second
property makes the lifted local component spherical once those characters are unramified.

% before: gal-solvable-local-base-change
For one cyclic step of prime degree, the second property comes from the local theory of Arthur and
Clozel. Their global lift has the local base change at every place, and local base change commutes with
parabolic induction.

% before: gal-solvable-base-change-facts
Each step of the tower must preserve the hypotheses of the next. Polarization, regular algebraicity and
weak attachment of the restricted Galois representation pass to a base change, and base changes along
two steps compose.

% before: gal-solvable-self-twist
Cyclic base change of prime degree is cuspidal unless the representation is isomorphic to a twist by a
character of the extension. Such a self-twist forces the trace of an attached Galois representation to
vanish off the subgroup fixing the extension, which an open subgroup with nonvanishing trace excludes.

% before: gal-solvable-cuspidal-lift
Induction along a composition series of the Galois group now carries a representation in scope along a
solvable CM extension disjoint from an open subgroup on which the trace does not vanish. At each step
Lemma~\ref{gal-solvable-self-twist} gives cuspidality and Lemma~\ref{gal-solvable-local-base-change} the
supports.

% before: gal-solvable-lift
The lift of Proposition~\ref{gal-plain-compatibility} is built from two such towers. After it, a given
element of nonzero degree fixes the top field, and the characters of the support are trivial on the
inertia elements that fix it.

% before: gal-plain-compatibility
At a place where the local component is plain, base change makes it spherical, and the spherical
conditions then determine the Weil--Deligne representation. The local correspondence on such a
component is a sum of characters by the laws, which gives full compatibility.

% before: gal-depth-zero-lift-layered
At a place over some $\mathbf Q_g$ with a local component of depth zero, the lift for
Proposition~\ref{gal-depth-zero-compatibility} passes to an abelian CM extension in which the place is
unramified with large residue degree. The extension is built one layer of prime degree at a time, and
each layer is ramified at a new place where the representation is unramified, so that cyclic base
change along it stays cuspidal. Such extensions are cut out by characters of the Galois group of the
totally real subfield that ramify only at split places.

% before: gal-depth-zero-lift-cyclotomic
The residue degree at the given place comes from a single character of finite order. It is cut out of
a cyclotomic field of prime conductor, with the prime chosen by the Chebotarev density theorem so that
the Frobenius element at the place has the required order.

% before: gal-depth-zero-lift-fresh
Each layer needs a new place of ramification at which the earlier layers are split. Characters
ramified at such a place again come from cyclotomic fields of prime conductor, with the prime split
completely in a given finite extension.

% before: gal-depth-zero-lift-field
The cyclotomic character has the right values at the given place, but its field need not be reached by
layers. Multiplying it by characters ramified at new places, one prime factor of its order at a time,
keeps those values and makes each step of prime degree ramified at a new place.

% before: gal-depth-zero-lift-weak
The representation is carried along the layers by cyclic base change. Only the Satake parameters at
almost all places are recorded along the way, so the base changes are weak in the sense defined here.

% before: gal-depth-zero-lift-base-change
Theorem~\ref{gal-cyclic-base-change} gives a lift when no twist of the required kind fixes the
representation. The construction also uses descent for Galois-stable representations and the fibres
of base change, which are part of the same theorem of Arthur and Clozel; it is stated here with weak
base changes and norm characters.

% before: gal-depth-zero-lift-step
A layer is ramified at a place where the representation is unramified, so no twist by a norm character
of the layer fixes the representation, and cyclic base change gives a cuspidal lift. The lift stays in
scope, and the restriction of an attached Galois representation stays attached.

% before: gal-depth-zero-lift-step-support
The proof above rules out a self-twist with the irreducibility of the global representation and an
averaging projection at the place $v$. The same conclusion follows from the local component at $v$ alone,
through the uniqueness of supercuspidal supports.

% before: gal-depth-zero-lift-tower
Iterating over the layers needs, at each layer, a place of the current field where the current
representation is unramified. The places split in the earlier layers serve, since a base change does
not change the local components at places of degree one.

% before: gal-depth-zero-lift-stable-spectrum
The stability of a weak base change under the Galois group of its step is proved in the space of
square-integrable functions on $\operatorname{GL}_n(K)\backslash\operatorname{GL}_n(\mathbf A_K)^1$. There
a cuspidal representation becomes an irreducible closed subspace, and its infinitesimal characters and
Satake parameters become eigenvalues of operators on that subspace. An automorphism of the field acts on
this space, and a stable subspace is one that it preserves.

% before: gal-depth-zero-lift-stable-parabolic
Base change acts on parameters by restricting embeddings and raising Satake parameters to the power of
the residue degree. The comparison of trace formulas produces a cuspidal summand with given parameters
only when these are not the parameters of the continuous spectrum, which come from cuspidal summands of
smaller general linear groups shifted by complex numbers. A weak base change to a sum allows the base
change of a cuspidal representation to be a sum of cuspidal representations.

% before: gal-depth-zero-lift-stable-carried
Cuspidal summands and cuspidal automorphic representations determine each other, after a twist by a
complex power of the norm of the determinant. Under this correspondence the parameters of a summand are
those of the representation it carries, and multiplicity one makes isomorphic representations equal.

% before: gal-depth-zero-lift-stable-twist
Composing a representation of $\operatorname{GL}_n(\mathbf A_M)$ with a field automorphism moves the
cuspidal summand that carries it by the Galois twist. So stability of the representation under a cyclic
Galois group is stability of its summand under the twist by a generator.

% before: gal-depth-zero-lift-stable-summand
Comparing the trace formula over $K$ with the trace formula over $M$ twisted by a generator of the
Galois group gives base change on the level of cuspidal summands. Its output is a summand stable under
the Galois twist with the base-changed parameters, provided that these avoid the parameters of parabolic
data over both fields.

% before: gal-depth-zero-lift-stable-summand-comparison
Alpöge proves the proposition with a comparison of trace formulas of his own, in which Hecke operators
remove the noncuspidal spectrum. Its inputs are matchings of test functions at every place and a
positivity argument over the smaller field.

% before: gal-depth-zero-lift-stable-avoid
The hypothesis of Proposition~\ref{gal-depth-zero-lift-stable-summand}(2) has to be checked for the
base-changed parameters of a representation that has a cuspidal weak base change. If a parabolic datum
over either field had these parameters, the cuspidal weak base change would have the Satake parameters
of a sum of at least two cuspidal representations, which strong multiplicity one rules out.

% before: gal-depth-zero-lift-infinity
The chain built along the layers records only Satake parameters. The infinitesimal characters are
recovered by descent: a weak base change is stable under the Galois group of its step, and two
representations with the same weak base change differ by a twist of finite order.

% before: gal-depth-zero-lift-infinity-stable
The proof above obtains the Galois stability of a weak base change from a stable summand produced by
the trace formula, and identifies its cuspidal members by strong multiplicity one. Applied to the
weak base change itself and its conjugates, strong multiplicity one gives the stability at once.

% before: gal-support-level-zero
The local component at the top of a chain of weak base changes is reached from the one at the bottom
through level-zero parts, which are representations of finite general linear groups. Green's cuspidal
characters and Harish-Chandra induction describe these parts, and Shintani lifting carries class
functions up an extension of residue fields.

% before: gal-support-finite
The following facts on finite general linear groups are classical. They identify the class functions
$\kappa_\Theta(m,c)$ with cuspidal characters, and they compute pairings with Harish-Chandra inductions
by Frobenius reciprocity.

% before: gal-support-pattern
A pattern records a Harish-Chandra induction of determinant characters and cuspidal characters that
the level-zero part of a representation meets. Lifting a pattern by a prime follows the effect of an
extension of residue fields of that degree on the parameters of the cuspidal characters.

% before: gal-support-pattern-arithmetic
Patterns are compared at a common residue field through their blocks, which list the characters of its
multiplicative group that the pattern produces. Lifting does not change the blocks, and once the residue
field is large enough for every cuspidal parameter, only characters remain.

% before: gal-shintani-situation
Shintani lifting compares class functions on general linear groups over a finite field and over an
extension of it, through the norm $\delta\mapsto\delta\,s(\delta)\cdots s^{l-1}(\delta)$ attached to a
generator $s$ of the Galois group. The next statements share the two groups, the norm and the block upper
triangular subgroups.

% before: gal-shintani-twisted
The Shintani lift of a Harish-Chandra induction is again an induction, in which conjugation is twisted by
the Galois generator. The lift of one cuspidal character is reached through the semidirect product of the
larger group with the Galois group, and through functions defined, like the cuspidal characters, by
induction over invariant subspaces.

% before: gal-shintani-norm
The facts on the norm map used below all come from Hilbert's Theorem 90 for the unit groups of Galois-stable
algebras of matrices. They show that every norm is conjugate to a matrix over the smaller field, move any
element by twisted conjugation into the algebra generated by its norm, and count the elements whose norms
lie in a given conjugacy class.

% before: gal-shintani-lift
Since twisted conjugation conjugates the norm, Shintani lifts are constant on twisted conjugacy classes. The
count of norms shows that lifting multiplies the inner product of class functions by the ratio of the orders
of the two groups.

% before: gal-shintani-parabolic
To compare the two inductions one must know which twisted conjugates of an element lie in a block upper
triangular group. For an element of the algebra generated by its norm, these come from matrices over the
smaller field that conjugate the norm into the block upper triangular group, multiplied by block upper
triangular matrices over the larger field.

% before: gal-shintani-hc
With the twisted induction, Shintani lifting commutes with Harish-Chandra induction. This reduces the lift
of the class function of a pattern to the lifts of its determinant characters and of its cuspidal
characters.

% before: gal-shintani-clifford
The support of the lift of a cuspidal character is read off from a virtual character of a semidirect product
by a group of prime order. Clifford theory for the subgroup of prime index turns two numerical facts into a
support statement: the restriction to a nontrivial coset has norm one, and the pairing on the subgroup with a
stable character is not divisible by the prime.

% before: gal-shintani-virtual
The virtual character is built by induction on the rank, from the tensor induction of a power sum of roots
of the characteristic polynomial and from inductions out of block upper triangular subgroups of the
semidirect product. The compatibility of Shintani lifting with Harish-Chandra induction identifies its
values on the coset of the Galois generator.

% before: gal-shintani-pairing
The pairing on the subgroup is computed from the pairings of cuspidal characters with power sums of roots.
When the prime divides the degree the pairing is a sign times the order of the group; otherwise it counts
tuples of exponents, and the count is $1$ modulo the prime.

% before: gal-shintani-cuspidal
For one cuspidal character these facts give the support of its Shintani lift. The lift is carried by the
class function of the pairs that replace the cuspidal pair when the pattern is lifted, which is the case of
Lemma~\ref{gal-support-shintani} with one cuspidal pair and no determinant blocks.

% before: gal-shintani-remark
Shintani's theorem describes the lifts of all irreducible characters at once. Together with the count of
norms, it gives the case of Lemma~\ref{gal-support-shintani} without determinant blocks directly.

% before: gal-support-shintani
At a place that is inert in a cyclic step, the level-zero parts below and above are compared through
Shintani lifting. The Shintani lift of the class function of a pattern acts by zero on every
representation that does not meet the class function of the lifted pattern.

% before: gal-support-depth-zero
The passage between representations of $p$-adic groups and their level-zero parts uses the theory of
depth-zero types and of types for principal series. The last two items turn a vector on which a
parahoric subgroup acts by a character into a supercuspidal support of characters, unlinked when the
characters on the blocks are distinct.

% before: gal-support-pattern-exists
A representation of depth zero has a pattern whose cuspidal entries are the level-zero characters of
the members of its support. Under the hypotheses of item (4) of Lemma~\ref{gal-depth-zero-lift-support},
the members of degree one can be grouped by their restrictions to the units, which gives an apart
pattern.

% before: gal-support-pattern-type
Once every cuspidal entry of a pattern has degree one, the pattern involves only determinant
characters. The level-zero part then contains a character of a parabolic subgroup, which gives a
vector of the kind used in Lemma~\ref{gal-support-depth-zero}(5).

% before: gal-support-matching
The comparison at an inert place rests on an identity of orbital integrals. The functions on
$\operatorname{GL}_n(\mathcal O)$ inflated from a class function on the residue field and from its
Shintani lift have matching orbital integrals under unramified base change.

% before: gal-support-stable
A weak base change along a cyclic step is stable under the Galois group of the step, and it does not
change the local components at places of degree one. Both facts are used at every step of the chain.

% before: gal-support-inert-trace
At an inert unramified place, the identity of traces of local base change, applied to the matching
functions, compares the level-zero parts below and above. A class function that pairs nonzero with the
level-zero part below has a Shintani lift that acts nonzero on the level-zero part above.

% before: gal-support-pattern-step
Stability, the identity of traces and Shintani lifting together carry a pattern through one cyclic
step. At a split place the local components agree; at an inert place the pattern is lifted by the
degree of the step.

% before: gal-support-pattern-chain
Induction along the chain carries a pattern from the bottom to the top. The hypothesis that inertia
fixes the top field keeps every place on the way unramified over the bottom.

% before: gal-support-inertia
The action of inertia is read through the tame reciprocity law. At a place over $g$ in which $g$ is a
uniformizer, the local Artin map on inertia is a residue of $t(\varpi)/\varpi$, and the laws give the
characteristic polynomials of the parameter of a depth-zero supercuspidal through its level-zero
character.

% before: gal-depth-zero-lift-support
At the given place the residue degree of the top field is divisible by the degree of every
supercuspidal in the support of the local component. There the lifted local component has a support of
tame characters, and the characteristic polynomials on inertia of the original parameter are products
of linear factors given by these characters.

% before: gal-depth-zero-lift-characters
For a representation whose support consists of characters, the laws give the characteristic
polynomials of its parameter on inertia. They are products of linear factors given by the characters
through the local Artin map, and they do not change under unramified twists.

% before: gal-depth-zero-lift
These pieces give the lift used in the proof of Proposition~\ref{gal-depth-zero-compatibility}. Over
the top field the local component has a plain twist, and the characteristic polynomials on inertia of
the two local parameters are given by the same tame characters.

% before: gal-depth-zero-compatibility
At places over some $\mathbf Q_g$ with local components of depth zero, a further base change reduces
to plain components. Only the restriction to inertia survives the descent, so the conclusion is about
characteristic polynomials on inertia. This is the form needed for the tame clause of compatibility.

% before: gal-row-from-core
The compatibility statements now give the full list of properties required of attached triples. The
two clauses of compatibility in Definition~\ref{gal-compatible} correspond to the plain and the
depth-zero cases. The Hodge--Tate multisets come from the infinitesimal character.

% before: gal-core-all
The core of the construction produces Galois representations only in the self-dual scope. Twisting by
Hecke characters removes the self-duality condition, and base change to CM extensions in good scope,
with gluing and descent, removes the splitting conditions.

% before: gal-definite-situation
The core starts from automorphic representations of definite unitary groups of odd rank. The fields
and places are chosen so that Theorem~\ref{gal-seed-representations} applies after a congruence. The
two places over $\ell$ are excluded from the conclusions.

% before: gal-generic-parameters
A seed has dimension larger than $a$, and the division step extracts an $a$-dimensional piece. It
identifies the Frobenius eigenvalues of that piece only where the Satake parameters admit no unexpected
multiplicative relation, which is the genericity defined here.

% before: gal-char-of-ann
At a place with multiplicatively generic Satake parameters, a matrix killed by the expected polynomial
and with the expected determinant has the expected characteristic polynomial. This is how the division
step determines Frobenius eigenvalues from the annihilating polynomials of the seeds.

% before: gal-sen
The division step needs an element of the Lie algebra of the image with $a$ distinct, well-separated
eigenvalues. Sen's theorem provides it from the lacunary Hodge--Tate weights of a seed.

% before: gal-lacunary-hull
The division is proved by linear algebra in $M_N(K)$, with the Zariski closure of the image and its
tangent space at the identity in place of an algebraic group and its Lie algebra. The objects defined
here recur throughout the proof: the envelope of the Lie algebra, systems of idempotents, and the
corner systems and matrix units that cut out an $n$-dimensional piece.

% before: gal-lacunary-specialize
The Lie element of Theorem~\ref{gal-sen} has entries in a large algebraically closed field, and the
division needs one with entries in $K$. A point over an integral domain can be moved to a point over
$K$ that satisfies the same polynomial equations and avoids finitely many hypersurfaces.

% before: gal-lacunary-descent
Membership in the Lie algebra and being an eigenvalue are polynomial conditions, and lacunarity is the
nonvanishing of finitely many linear forms. Specialization therefore brings the Lie element and its
lacunary eigenvalues down to $K$.

% before: gal-lacunary-closure
The proof uses a few facts about Zariski closures of monoids of matrices: they are closed under
products, the tangent space acts on their equations by derivations and is closed under brackets, its
exponential satisfies their equations, and a degree bound on the monoid passes to both. These are the
familiar facts about algebraic groups and their Lie algebras, proved here from the definitions.

% before: gal-lacunary-finite-index
The proof builds elements of the Zariski closure of the image of a subgroup of finite index from the
Lie algebra of the whole image. This is possible because the Lie algebra of the whole image lies in
the tangent space of the image of every normal subgroup of finite index.

% before: gal-lacunary-envelope
The associative algebra generated by the Lie algebra carries the structure used below. Conjugation by
the image preserves it, and it preserves every invariant subspace, as the Lie algebra does.

% before: gal-lacunary-semisimple
Matrix units are built in an algebra in which no nonzero element generates a nil left ideal. For an
algebra normalized by a semisimple representation and preserving its invariant subspaces, this follows
from the existence of invariant complements.

% before: gal-lacunary-grading
A lacunary element of the Lie algebra with $n$ eigenvalues decomposes $K^N$ into $n$ eigenspaces.
Lacunarity makes the differences of its eigenvalues distinct, so the Lie algebra splits into the
corresponding blocks.

% before: gal-lacunary-cycles
The degree bound $n$ allows at most $n$ distinct eigenvalues. An invertible element that permutes the
idempotents along the cycles of a permutation already has $n$ eigenvalues when distinct cycles share
none, and the bound then forces it to be scalar on each idempotent that the permutation fixes.

% before: gal-lacunary-diagonal
An element of the Lie algebra that commutes with the idempotents can be added, with any coefficient,
to the lacunary element without leaving the Lie algebra, so the degree bound holds for the sum. The
previous lemma, applied to the sum for a generic coefficient, shows that the element is scalar on each
block.

% before: gal-lacunary-algebra
The diagonal blocks of elements of the Lie algebra are scalar, and its off-diagonal blocks lie in it.
Both properties pass to the associative algebra it generates, because each product of blocks is a
commutator or a scalar multiple of a block.

% before: gal-lacunary-structure
The previous lemmas combine into a description of the envelope relative to the idempotents of the
lacunary element. Each idempotent cuts a one-dimensional corner out of the envelope, which is what the
construction of matrix units needs.

% before: gal-matrix-units
In a semisimple algebra with one-dimensional corners, a nonzero element between two corners has an
inverse between the same corners. This gives matrix units linking all idempotents in the same class.

% before: gal-matrix-units-conjugate
The image of $\rho(g)$ need not lie in the envelope, and its effect on the corners is controlled by
comparing the idempotents with their conjugates by $\rho(g)$. In a semisimple algebra two corner
systems of the same size are conjugate by a unit, up to a permutation of the indices.

% before: gal-matrix-units-factor
A unit of the envelope is a product of unipotent elements built from off-diagonal blocks and an element
that is diagonal in the idempotents. The unipotent factors lie in the Zariski closure of the image,
which brings the comparison of corner systems back into the image of $\rho$.

% before: gal-lacunary-generic
After the conjugation, $\rho(g)$ times a unit permutes the idempotents. Multiplying it by a diagonal
element whose entries satisfy no multiplicative relation from a finite list separates the eigenvalues
belonging to different cycles.

% before: gal-lacunary-torus
The diagonal multipliers must lie in the Zariski closure of the image. The exponential of the lacunary
element provides a point over $K[[t]]$ with the required independence, and specialization gives one
over $K$.

% before: gal-lacunary-key
The previous lemmas combine into two properties of each element of the group. For every coset of a
normal subgroup of finite index, $\rho(g)$ times some element of the Zariski closure of the image of
the subgroup has $n$ distinct eigenvalues, and on each corner that it meets $\rho(g)$ agrees with an
element of the envelope.

% before: gal-lacunary-big-corner
The envelope need not contain the image of $\rho$. Adjoining the image keeps the corners
one-dimensional, by Proposition~\ref{gal-lacunary-key}(2), and gives an algebra in which matrix units
define a representation.

% before: gal-matrix-units-representation
The coefficients of $\rho(g)$ with respect to the matrix units form an $n\times n$ matrix, and the
multiplication rules of the matrix units make this a homomorphism. Columns built from the matrix units
realize it inside $\rho$.

% before: gal-matrix-units-subrepresentation
The representation $r$ sits inside $\rho$, and $\rho$ is generated by copies of it. For semisimple
$\rho$ this makes $r$ semisimple, gives both the same annihilating polynomials, and splits $\rho$ off a
sum of copies of $r$.

% before: gal-lacunary-coset
The density of elements with separable characteristic polynomial is reduced to cosets of subgroups of
finite index. A polynomial condition that holds on an open subset of a compact group holds on a whole
coset of such a subgroup, by an argument with irreducible components.

% before: gal-lacunary-separable
Separability of the characteristic polynomial is a polynomial condition. An element of the Zariski
closure with $n$ distinct eigenvalues forces it at some element of the corresponding coset, and the
previous lemma turns this into openness and density.

% before: gal-division
The division is a statement about representations of a compact group: if a dense set of elements is
killed by polynomials of degree $n$ with a consistent constant term, and the Lie algebra contains a
lacunary element, then the representation is built from a single $n$-dimensional piece. Applied to a
seed, it produces the $a$-dimensional representation.

% before: gal-division-seed
Applied to the seeds of Theorem~\ref{gal-seed-representations}, with the determinant given by the
central character of $U$, the division produces a representation of dimension $a$. Its local properties
are inherited from the seed because it is a subrepresentation of the seed and the seed is a quotient of
copies of it.

% before: gal-congruences-calculus
The proof of Proposition~\ref{gal-congruences} approximates systems of Hecke eigenvalues in several
stages. Two formal properties let the stages be chained: approximations compose, keeping a polynomial
condition that the intermediate systems satisfy, and an approximation of a family restricts to any
subfamily of its coordinates.

% before: gal-congruences-eigen
Each stage obtains its approximation from one estimate between eigenvectors. If a normalized eigenvector
is $\ell$-adically close to a sum of other eigenvectors, measured by evaluations that the operators do not
enlarge, then every integral polynomial in its eigenvalues is close to a fixed combination of the same
polynomial in theirs.

% before: gal-congruences-forms
The congruences are found among algebraic modular forms on the definite group $U_a$, which take values in
the representation of the weight at infinity. The forms whose values lie in automorphic representations
with a given property span a subspace, and each stage compares a form with a combination of forms in such a
subspace.

% before: gal-congruences-dense
All stages produce the same kind of estimate: a form is close to a combination of forms of the desired
kind, with a relative error that is $\ell$-adically small at every point. To turn this into estimates for
eigenvalues, the Hecke operators must not enlarge the $\ell$-adic size of the forms, which holds for
integral operators measured on a gauge.

% before: gal-congruences-gauge
The gauge twists the argument of the representation at infinity by the components at the places over
$\ell$, which turns the action at infinity into an $\ell$-adic action. Integral Hecke operators, scaled by
the denominators that this twist introduces, then preserve the bound on the evaluations.

% before: gal-congruences-density
Density of a property in the space of forms implies the approximation of eigenvalue systems. A normalized
eigenform with values in $U$ is close to a combination of eigenforms with the property, and
Lemma~\ref{gal-congruences-eigen} converts this into an approximation of its eigenvalues.

% before: gal-congruences-weight-family
To change the weight at $v'$, a form of weight $\lambda$ is multiplied by functions of the $\ell$-adic
component at $v'$ that behave at infinity like a highest weight vector of weight $\nu$. The resulting family
stays $\ell$-adically close to the form, and its automorphic components have weights within a fixed
distance of $\nu$ at the embedding over $v'$.

% before: gal-congruences-weights
Choosing $\nu$ with entries that grow geometrically forces the weights of the components to be lacunary at
the embedding over $v'$. The Hecke operators away from $v'$ commute with the twist, so the functions built
from the family are eigenvectors to which Lemma~\ref{gal-congruences-eigen} applies.

% before: gal-congruences-section
The type at $w_0$ and the genericity are both obtained by multiplying a form by functions of the class of
its argument modulo a smaller level. Such functions come from a homomorphism from a level to a finite group,
extended to the whole group compatibly with the left action of $U_a(F)$, which is possible when the
homomorphism kills the elements of finite order.

% before: gal-congruences-type-character
The supercuspidal type at $w_0$ comes from a character of a small congruence subgroup. It is defined
through the trace against a matrix with no invariant subspace modulo the maximal ideal, which makes it
nontrivial on the unipotent radical of every proper parabolic subgroup and so kills the Jacquet modules.

% before: gal-congruences-type
A form fixed by the congruence subgroup $J$ at $w_0$ is congruent modulo $\ell^m$ to a combination of its
twists by the nontrivial characters of $\mathbf Z/\ell^m$. The congruence comes from the sum of all
characters, which is $\ell^m$ at $0$ and vanishes elsewhere.

% before: gal-congruences-l-group
Genericity rests on an algebraic fact about $\ell$-groups. A translation-invariant space of functions on a
finite $\ell$-group whose dimension is at least half the order of the group contains a function congruent
to $1$ modulo $\sqrt\ell$.

% before: gal-congruences-negligible
At a split place the Satake isomorphism identifies spherical Hecke operators with symmetric Laurent
polynomials in the Satake parameters. A polynomial equal to $1$ on the zero set of $\Phi$ and small in mean
square on the unitary torus gives, by the Plancherel formula, an operator of small norm that acts as the
identity on the non-generic representations.

% before: gal-congruences-sparse
An operator of small norm bounds the dimension of any space of forms on which it acts as the identity. In a
deep enough level of the tower of congruence subgroups at $w_0$, the non-generic forms therefore have
dimension at most half the index of that level in the first one, which is the bound that
Lemma~\ref{gal-congruences-l-group} needs.

% before: gal-congruences-generic
Together these give the density of generic forms. Multiplying a form by the functions on the $\ell$-group
$J_0/J_n$ produces a space of which the generic part has at least half the dimension, so it contains a
multiple of the form congruent to it modulo $\sqrt\ell$; iterating improves the congruence.

% before: gal-congruences-operators
The proof of Proposition~\ref{gal-congruences} tests genericity, power sums of Satake parameters and the
given Hecke eigenvalues through global Hecke operators. They come from local operators at split places
through the Satake isomorphism, and a power of $\ell$ makes each of them $\ell$-integral.

% before: gal-congruences
The division step needs a lacunary weight, the supercuspidal type at $w_0$ and generic Satake
parameters, which $U$ need not have. Congruences between algebraic modular forms on the definite group
replace $U$ by forms with these properties whose Hecke eigenvalues approximate those of $U$
$\ell$-adically.

% before: gal-trace-limit
The representations attached to the approximating forms are combined by passing to the limit in their
traces. Here approximation means uniform approximation of all polynomial expressions by linear
combinations, and the limit of traces is the trace of a representation by Taylor's theorem on
pseudo-characters.

% before: gal-cover-definition
Alpöge passes the Weil--Deligne traces at places over $\ell$ to the limit of
Lemma~\ref{gal-closed-crystalline}(2) by comparing integral crystalline periods of $r$ and of the
approximating representations. The comparison takes place on lattices: one lattice for $r$, and at
each level a lattice for finitely many approximating representations, related to the first by a matrix
that is equivariant modulo $\ell^n$. Weighted by the coefficients of the approximation, the traces on
the two lattices agree up to $\ell^{-n}$.

% before: gal-cover-joint
A cover is built from approximations of the trace of $r$ at finer and finer precision. When the same
coefficients also approximate a second number $\theta$, the layer at level $n$ approximates $\theta$
to within $\ell^{-n}$, and this is the form in which the limit is taken. The semisimplicity of $r$
bounds the denominators that appear in the construction.

% before: gal-cover-breuil-kisin
Crystalline Frobenius operators on different lattices are compared through Breuil--Kisin modules. A
Breuil--Kisin structure on a lattice records the Frobenius of a Breuil--Kisin module, its comparison
with the lattice over $W(\mathbf C_\ell^\flat)$, and the Galois action over $A_{\inf}$. A congruence
between two lattices becomes a congruence between their structures, called a link.

% before: gal-structures-datum
A Breuil--Kisin datum presents $\mathcal O_K$ as the quotient of a power series ring over Witt vectors by
an Eisenstein polynomial, and maps that ring to $A_{\inf}$ compatibly with Frobenius and with Fontaine's map
$\theta$. Every finite extension of $\mathbf Q_\ell$ has one, built from a uniformizer and the unramified
subfield.

% before: gal-structures-twist
A cyclotomic element is obtained from Fontaine's generator of $\ker\theta$ by correcting a unit. The
correction solves the equation $\phi(\eta)=\omega\eta$ in the units of $A_{\inf}$, which is possible because
every monic polynomial over $\mathcal O^\flat$ has a root there.

% before: gal-structures-cyclotomic
The cyclotomic element plays the part of a period of the cyclotomic character for the Breuil--Kisin datum.
Its Frobenius is $\iota(E)$ times itself, and it is invertible in $W(\mathbf C_\ell^\flat)$, so comparison
matrices can be multiplied and divided by its powers.

% before: gal-structures-integral
Integrality over $A_{\inf}$ can be read off from a Frobenius equation. Frobenius multiplies a negative
valuation by $\ell$, so a matrix whose Frobenius is an integral linear expression in the matrix itself
cannot have entries with poles.

% before: gal-structures-basis
The Galois action over $A_{\inf}$ of a Breuil--Kisin structure is determined by its comparison matrix. On
the lattice of a cover the matrices $T$, $A$ and $A'$ are constructed first, and the action is checked to
be integral afterwards.

% before: gal-structures-poles
Multiplying the comparison matrix of a Breuil--Kisin basis of height $N$ by $\tau^N$ makes it and its
inverse integral. This turns the Frobenius equation of the basis into one with integral coefficients, to
which Lemma~\ref{gal-structures-integral} applies.

% before: gal-structures-saturated
Congruences modulo all powers of $\ell$ detect integrality over $A_{\inf}$, even after multiplication by a
fixed power of $\ell$. Integrality and the crystalline condition pass in this way from the layers of a
cover to its lattice.

% before: gal-structures-layer
The representations of the layers of a cover are built from members of $C$ by restriction of scalars and
direct sums. They inherit crystallinity and the bound on Hodge--Tate weights, which are what the
existence of Breuil--Kisin structures requires.

% before: gal-structures-kisin
Lattices in crystalline representations with Hodge--Tate weights in $[-N,N]$ carry crystalline
Breuil--Kisin structures of height $N$. This is Kisin's theory of Breuil--Kisin modules, with the Galois
action over $A_{\inf}$ from the theories of Liu and Gao and the crystalline condition from the work of Gee, Liu
and Savitt and, for every $\ell$, of Du and Liu.

% before: gal-structures-functorial
An equivariant map between lattices with Breuil--Kisin bases is given in those bases by a matrix over
$S$. This is the full faithfulness of Kisin's functor, and it supplies the transition matrices along a
lattice tower.

% before: gal-structures-etale
The Breuil--Kisin basis of the lattice of a cover is approximated first by bases over a ring of Laurent series
in $u$, completed $\ell$-adically. This ring sits inside $W(\mathbf C_\ell^\flat)$ and is cut out by approximation
modulo powers of $\ell$.

% before: gal-structures-topological
Coefficients over $S$ are read off from elements of $A_{\inf}$ through a topological basis of $A_{\inf}$ over $S$.
Its coordinate maps are $S$-linear, so they carry relations over $A_{\inf}$ to relations over $S$.

% before: gal-structures-topological-exists
A topological basis exists because $\ell$ and $u$ form a regular sequence on $A_{\inf}$, for whose ideal $A_{\inf}$
is complete. A basis of $A_{\inf}$ modulo $(\ell,u)$ over the residue field of $S$ then lifts by successive
approximation.

% before: gal-structures-frobenius-unit
A power series not divisible by $\ell$ whose Frobenius divides it up to a power of $E$ is a unit. The proof
compares absolute values of roots in $\bar{\mathbf Q}_\ell$, and it uses that $E$ is not a polynomial in
$u^\ell$.

% before: gal-structures-coefficients
The coordinates of a power of the cyclotomic element generate an ideal of $S$ that contains a power of $\ell$.
This bounds the defect of the right inverses in the limit construction, and it recovers coefficients over $S$
from multiples that are integral over $A_{\inf}$.

% before: gal-structures-integral-etale
An element of Fontaine's ring that lies in $A_{\inf}$ comes from $S$. The proof removes one power of $\ell$ at a
time by comparing zeroth Witt components.

% before: gal-structures-descent
Matrices over Fontaine's ring descend to $S$ in two ways: by integrality after multiplication by a power of the
cyclotomic element, and by approximating solutions over $A_{\inf}$ modulo powers of $\ell$. Both read coefficients
through a topological basis.

% before: gal-structures-free
The limit Breuil--Kisin basis spans a module over $S$ cut out of $A_{\inf}^d$ by a condition over Fontaine's ring.
Such a module is free of rank $d$, because modulo $\ell$ it becomes a lattice over the principal ideal domain
$\mathbf F_{\ell^h}[[u]]$.

% before: gal-structures-points
The bases over Fontaine's ring used for the lattice of a cover come from showing that certain elements of
$\mathbf C_\ell^\flat$ lie in the reduction of that ring. This is tested through maps to perfect valuation rings
that are complete for the image of a given element; they serve as points of the rings involved.

% before: gal-structures-points-bounded
Divisibility of the powers of an element by powers of $u$, up to a fixed defect, is tested on valuation rings.
Either the defect is bounded, or a valuation ring of an algebraically closed field sees the element as dividing
$u^M$.

% before: gal-structures-points-completion
The valuation rings found in this way need not be complete. Completing along the image of $u$ keeps them
valuation rings with bijective Frobenius, so they become perfect points.

% before: gal-structures-points-killed
An element that satisfies a Frobenius equation and is killed by every perfect point is divisible by every power
of $u$. The proof inverts the elements $1-ur$, bounds the powers of the element with perfect points, and iterates
the Frobenius equation.

% before: gal-structures-points-descent
Elements of a flat algebra on which any two maps agreeing on the base coincide descend to the base. The agreement
is required only at perfect points, together with a Frobenius equation for the differences, and it is applied to
the two inclusions into the tensor square.

% before: gal-structures-points-untilt
A perfect point $O$ carries its Witt vectors $W(O)$, and an element of the shape of the image of the Eisenstein
polynomial generates a prime ideal there. The quotient is the ring in which Fontaine's map $\theta$ is compared
for two perfect points.

% before: gal-structures-residue
Modulo $\ell$, Fontaine's ring becomes a field of Laurent series in $\pi^\flat$ inside $\mathbf C_\ell^\flat$. Its
ring of power series $\mathbf F_{\ell^h}[[u]]$ is a retract of $\mathcal O^\flat$ through a topological basis of
$A_{\inf}$, and the retraction is the map to the base in the descent along $\mathcal O^\flat$.

% before: gal-structures-points-galois
Two perfect points of $\mathcal O^\flat$ that agree on the power series in $\pi^\flat$ differ by an element of
$G_K$ fixing $\pi^\flat$. The element comes from comparing two maps of $\mathcal O_{\mathbf C_\ell}$ into the
quotient of $W(O)$ by the image of the Eisenstein polynomial, which agree on $\mathcal O_K$.

% before: gal-structures-residue-descent
An element of $\mathbf C_\ell^\flat$ fixed by $G_\infty$ lies in the reduction of Fontaine's ring if any two
copies of it satisfy a Frobenius equation over that reduction. The descent along perfect points applies, and the
Galois element of Lemma~\ref{gal-structures-points-galois} gives the agreement at points.

% before: gal-structures-echelon
The kernel of a matrix with a Frobenius equation has a projector in echelon form. The projector inherits the
Frobenius equation and every Galois invariance of the kernel, so its entries satisfy the hypotheses of the
descent along perfect points.

% before: gal-structures-kernel-vector
A singular matrix over $\mathbf C_\ell^\flat$ with a Frobenius equation over the reduction of Fontaine's ring,
Galois invariant up to invertible matrices, has a kernel vector over that reduction. The vector is a column of
its echelon projector.

% before: gal-structures-unit
A matrix over $W(\mathbf C_\ell^\flat)$ that is stable under Frobenius and under $G_\infty$ is invertible if no
vector over Fontaine's ring outside $\ell$ times that ring is killed modulo $\ell$. Modulo $\ell$, a singular
such matrix would have a kernel vector over the reduction of Fontaine's ring.

% before: gal-structures-saturated-lattice
Multiplying a matrix on the right by a basis of the vectors over Fontaine's ring that it sends to multiples of
$\ell^c$, for the power $\ell^c$ in its determinant, and dividing by $\ell^c$, gives a matrix that kills modulo
$\ell$ only vectors divisible by $\ell$. This saturation is what the invertibility criterion asks for.

% before: gal-structures-isogenous
A Galois-stable sublattice of finite index in a lattice with a Breuil--Kisin basis has a basis over
$W(\mathbf C_\ell^\flat)$ that differs from the Breuil--Kisin basis by a matrix over Fontaine's ring. The
saturated matrix of the sublattice is stable under Frobenius and under $G_\infty$, so it is invertible.

% before: gal-structures-kernel-lattice
The vectors that a congruence $\Pi$ sends to zero modulo $\ell^n$ form a sublattice stable under the Galois
action. Its basis matrix intertwines the action on the sublattice with the original one.

% before: gal-structures-smith
Elementary divisors over Fontaine's ring factor a matrix with a right inverse modulo $\ell^n$ as an invertible
matrix times a matrix over Fontaine's ring, when the vectors it kills modulo $\ell^n$ are spanned by a matrix over
that ring. The columns that survive modulo $\ell^n$ are counted over the residue field.

% before: gal-structures-etale-basis
The limit construction of Lemma~\ref{gal-structures-limit} needs, for each lattice of the tower, a basis over
$W(\mathbf C_\ell^\flat)$ in which the image of its Breuil--Kisin basis has coefficients in Fontaine's ring
modulo $\ell^n$. It is assembled from the kernel lattice of $\Pi$, its basis over Fontaine's ring and the
elementary divisor factorization.

% before: gal-structures-limit
The lattice of a cover is a limit of the layer lattices only modulo powers of $\ell$, so its Breuil--Kisin
basis is constructed as a limit. A basis over Fontaine's ring $\mathcal O_{\mathcal E}$ is built first, and the
Breuil--Kisin basis is cut out of it by integrality over $A_{\inf}$.

% before: gal-structures-limit-frame
A Breuil--Kisin basis linked to Breuil--Kisin structures with bounded defect inherits their Galois action.
The action is integral because, up to a fixed power of $\ell$, it is congruent to integral actions modulo
every power of $\ell$, and the crystalline condition passes in the same way.

% before: gal-cover-structures
The lattices of a cover carry Breuil--Kisin structures of a single height, since the Hodge--Tate weights
of the approximating representations are fixed. The structure on the lattice of $r$ is linked to those
of the layers with arbitrary precision, and multiplying its comparison matrix by a power of the
cyclotomic element makes it integral.

% before: gal-cover-weil-frobenius
In Fontaine's construction an element of the Weil group of degree $d$ acts on crystalline periods
through $\phi^{-fd}$. The powers of $\phi$ that occur must fix the constants of the Breuil--Kisin datum,
which lie in the maximal unramified subfield of $K$.

% before: gal-cover-zero
Integral identities in $A_{\mathrm{cris}}$ are proved by showing that a defect is divisible by $u$ and
satisfies a Frobenius equation. Frobenius raises the power of $u$, and the divided powers of
$A_{\mathrm{cris}}$ turn high powers of $u$ into high powers of $\ell$, so the defect vanishes or is
small modulo a power of $\ell$.

% before: gal-cover-section
The Frobenius matrix of a Breuil--Kisin module and its constant term are conjugate over
$A_{\mathrm{cris}}$ up to a power of $\ell$. The conjugating matrix is a limit of iterated Frobenius
twists, and the power of $\ell$ depends only on the height, so the same power serves every lattice of a
cover.

% before: gal-cover-period
A crystalline Breuil--Kisin structure and the conjugating matrix give a period matrix of the lattice
over $B_{\mathrm{cris}}$ with controlled denominators. Relating two Frobenius twists of the period
matrix, a matrix $V$ over $A_{\mathrm{cris}}$ then represents the action of a Weil group element on
crystalline periods, up to a scalar $s$ that depends only on the height and the twists.

% before: gal-cover-link
Linked Breuil--Kisin structures give congruent integral period matrices. The link passes to $A_{\inf}$
and to the Frobenius matrices of the two modules, and the conjugating matrices of the two structures are
then congruent by Lemma~\ref{gal-cover-zero}.

% before: gal-cover-operators
Along a cover the previous constructions fit together. One scalar $s$ serves the lattice of $r$ and
every layer, and the matrices $V$ of the layers are congruent to that of $r$ along the lattice tower to
arbitrary precision; through these congruences the Weil--Deligne traces of the layers converge.

% before: gal-cover-trace-identity
The matrix $V$ computes Weil--Deligne traces. If a lattice reads $r$ through matrices $b_i$ and a vector
$e$, then pairing $Ve$ with the traces of the $b_i$ gives $s$ times the Weil--Deligne trace of $r$, after
any $\mathbf Q_\ell^{\mathrm{ur}}$-linear functional is applied to the traces.

% before: gal-cover-functional
The trace identities hold only after a $\mathbf Q_\ell^{\mathrm{ur}}$-linear functional is applied to
elements of $\bar{\mathbf Q}_\ell$. To detect that a difference of traces is nonzero, the functional must
not shrink it while keeping finitely many other elements integral.

% before: gal-cover-functional-free
The proof above builds the independent family one element at a time, by exchanging a member for the new
element when its coefficients are too large. Over a valuation ring the family also comes at once from the
structure of finitely generated torsion-free modules.

% before: gal-cover-trace-limit
The trace identity for the lattice of $r$ is compared with those for the layers. The congruences of the
matrices $V$ and the trace estimates of the cover make $s$ times the difference of the two Weil--Deligne
traces divisible by every power of $\ell$ in $A_{\mathrm{cris}}$, and $s$ is a unit of
$B_{\mathrm{cris}}$.

% before: gal-closed-crystalline
The local conditions at places over $\ell$ must survive the limit. Crystallinity with fixed Hodge--Tate
weights is closed under such limits, and so are the power sums of the crystalline Frobenius when they
are approximated jointly with the traces.

% before: gal-limits
Seeds exist only for automorphic representations with a supercuspidal type and lacunary weights. A
general $U$ is reached by $\ell$-adic congruences with such representations, and its Galois
representation is a limit of theirs. The conditions at places over $\ell$ survive the limit because
the Hodge--Tate weights are fixed.

% before: gal-patch-index-two
Stage 4 of the proof of Proposition~\ref{gal-core} builds representations of $G_Q$ for infinitely many
quadratic extensions $Q$ of a CM field and needs one representation of the Galois group of the base
field that restricts to them. The groups $G_Q$ are the kernels of distinct quadratic characters, and
semisimple representations are determined by their traces. Such a family can be glued at the cost of
finitely many members.

% before: gal-patch-unique
The trace conditions of Lemma~\ref{gal-patch-index-two} will come from Frobenius eigenvalues. Places of
prime norm suffice for this by the Chebotarev density theorem, and their norms do not change in a
finite extension.

% before: gal-patch-system
The representations to be glued are tied together by Satake parameters of the base field, read at the
places of prime norm of each extension. A place of prime norm has the same norm as the place of the base
field under it, so the eigenvalues $q_w^{(a-1)/2}B$ make sense over every finite extension.

% before: gal-patch-quadratic
Uniqueness turns the Frobenius data into the two trace conditions of Lemma~\ref{gal-patch-index-two}.
Conjugation by $G_F$ and restriction to the compositum of two members both preserve the property of
realizing a Satake system, so each representation is invariant under $G_F$ and any two agree where both
are defined.

% before: gal-patch-serves
To glue representations attached to a cuspidal $\pi$, each quadratic extension $Q$ must carry a
representation for the base change of $\pi$, with the four conditions at enough places. The conditions
at a place $w_0$ of the base field can be read over $Q$ when $w_0$ splits in $Q$. A ramified place over
a spherical place of $\pi$ rules out that $\pi$ is a twist of itself by the character of $Q$, so the
base change exists.

% before: gal-patch-descent
With these properties the gluing descends the four conditions from quadratic extensions to the base
field. Each place needs its own infinite family of extensions split there, and one glued representation
serves all places, since only finitely many extensions are lost in the gluing.

% before: gal-patch-steps
Stage 4 reaches an arbitrary CM field through three quadratic extensions. Two real steps split the
places over $\ell$ until every CM extension has three of them over distinct places of its maximal
totally real subfield, as admissibility requires. A CM step then makes the field inert-admissible while
the places of a given finite set stay split.

% before: gal-patch-steps-kummer
Both kinds of step are made by adjoining the square root of an integer, so their splitting and
ramification come from the action of decomposition and inertia groups on such a root. For a prime that
divides the integer once at a place, inertia acts on its $m$-th roots through all $m$-th roots of unity
when $m$ is prime to it. The proof is tame Kummer theory: an inertia-fixed power of the root would lie in
a field unramified at the place.

% before: gal-patch-steps-squares
For square roots this gives three tests. A congruence modulo $8t$ fixes a square root at the places over
the prime $t$, and a unit fixes it under inertia. By the previous lemma a prime that divides the integer
once, at a place unramified over that prime, moves it under inertia.

% before: gal-patch-steps-squares-inertia
For square roots the action of inertia can also be read off from valuations. The fixed field of inertia
is unramified, so it cannot contain a square root of a prime of valuation one.

% before: gal-patch-steps-frobenius
Where a finite set of elements is fixed by inertia, the decomposition group acts on it through the
powers of one Frobenius element. Questions about decomposition groups then become questions about the
parity of an exponent.

% before: gal-patch-steps-adjoin
Each step must again be a CM field. Under every complex embedding the square root of an integer is real
or purely imaginary, so complex conjugation extends from a CM field to the field obtained by adjoining
it.

% before: gal-patch-steps-cm-places
The conditions of a CM step concern the places not split over the maximal totally real subfield.
Splitting over that subfield is read from the decomposition group. At a place inert in a CM field, an
unramified square root of an integer is fixed by the decomposition group over the CM field, since that
group acts through even powers of Frobenius.

% before: gal-patch-steps-chebotarev
The CM steps adjoin the square root of minus a product of two primes. At the places over these primes,
every element of the decomposition group that fixes the given real quadratic element must fix the whole
field. The Chebotarev density theorem gives one such prime for each automorphism of the right shape,
with its action on a root of unity prescribed.

% before: gal-patch-steps-two-primes
Two such primes, for complex conjugation composed with an automorphism that negates $\delta$ and for
the inverse of that automorphism, have a product congruent to $-1$ modulo any given integer. Minus
their product is then a square locally at the places of a given finite set.

% before: gal-patch-steps-exist
Lemma~\ref{gal-patch-descent} needs infinitely many extensions of each kind for every finite set of
places. The CM steps are constructed over a field obtained by adjoining the square root of a positive
integer.

% before: gal-larger-enlargement
The construction of Stage 1 needs odd rank. For even $n$, a theta lift carries an automorphic
representation of $U_n$ to one of $U_{n+t}$ with $n+t$ odd, adding $t$ Satake parameters and $t$
Hodge--Tate weights that come from algebraic Hecke characters, after a twist by one more Hecke
character. Two such lifts whose added characters all differ determine the common $n$-dimensional part.

% before: gal-larger-good-base
Theorem~\ref{gal-theta-lift-large-rank} needs a good base and Hermitian and skew-Hermitian spaces with
self-dual lattices. On a CM field unramified over its maximal totally real subfield, the skew-Hermitian
lattice comes from the different, once global class field theory writes a suitable idele of the real
subfield as a global element times a norm.

% before: gal-larger-conjugate-character
The characters $\chi_X$ and $\chi_W$ of the theta lift must restrict to powers of $\eta$ with prescribed
infinity types. They are built as powers of one character with odd infinity type at every infinite
place.

% before: gal-larger-attuned
After the lift, the twist and the added characters are read off from $\chi_X$, $\chi_W$ and the norm.
Their infinity types and their values at uniformizers must satisfy linear relations for the weights and
Satake parameters of the lift to have the shape of an enlargement.

% before: gal-larger-characters
Two theta lifts are needed, with added characters that differ. Taking two different even powers of the
same character for $\chi_X$ separates the added characters, while the twist $\kappa$ is the same for both
lifts.

% before: gal-satake-support-gap
The local theta lift at a split place is an induced representation whose irreducibility rests on a bound
for the supercuspidal support of a unitarizable representation. The proof of that bound singles out, among
the characters of the support that differ from a given one by integral powers of the absolute value, a
block separated from the others by a gap of at least $2$. Finding such a block is a statement about finite
sets of integers.

% before: gal-satake-support-front
Suppose that no member of one part of a supercuspidal support is a twist by $\nu^{-1}$, $1$ or $\nu$ of a
member of the rest. The members of that part can then be moved to the front of a product of supercuspidal representations without
changing the product, so an irreducible representation with this support embeds into a product whose first
factor is built from that part alone.

% before: gal-satake-unitarizable-exponent
A unitarizable admissible representation has bounded matrix coefficients. For an embedding into a
parabolically induced representation, this bounds the central character of the first inducing factor at
the scalars $z$ with $|z|<1$.

% before: gal-satake-unitarizable-support
For $r<t$ and $\chi$ unitary, the character $\chi|\cdot|^{-(t-1)/2}$ cannot occur in the supercuspidal
support of a unitarizable representation of $\operatorname{GL}_r$. Otherwise the separated block containing
it could be moved to the front of an embedding, and the sum of its exponents would be too negative for the
bound on central characters.

% before: gal-satake-weil-equivariant
The second step of the local theta lift analyzes a map out of the Weil representation of a pair of general
linear groups according to the rank of matrices. The analysis needs only the action of
$\operatorname{GL}_r\times\operatorname{GL}_N$ on Schwartz--Bruhat functions on $N\times r$ matrices,
twisted by a character, and maps equivariant for it.

% before: gal-satake-weil-rank-ball
On matrices of rank at most $j$, the indicator function of a small ball around $e_j$ is the sum of
$q^{j(r-j)}$ of its translates under $\operatorname{GL}_r$. When $2j<N$ the normalizing factor of the Weil
representation exceeds the number of translates, and unitarity forces an equivariant map that already
vanishes from rank $j+1$ to vanish on that indicator function.

% before: gal-satake-weil-rank-step
Translating the ball by the group and localizing on the support extend the vanishing from rank $j+1$ to rank
$j$. Repeating the step descends through all the ranks below $r$.

% before: gal-satake-weil-germ
On the open set of matrices of rank $r$ an equivariant map is determined by its germ at $e_r$, a linear map
from the space of $\pi$ to that of $\sigma$. The germ vanishes only if the map vanishes from rank $r$, and under the parabolic
subgroup stabilizing the column space of $e_r$ it transforms by the inducing character.

% before: gal-satake-weil-lift
If the germ were zero, the rank steps would make the equivariant map vanish everywhere. So the germ is
nonzero, and by Frobenius reciprocity it gives a nonzero map from $\pi$ to the induced representation, which
is an isomorphism when the induced representation is irreducible.

% before: gal-larger-split-satake
At a split place the definite unitary groups become general linear groups, and the local theta lift is
a correspondence for the pair $(\operatorname{GL}_n,\operatorname{GL}_d)$. In the range $d>2n$ the
contragredient of the lift of a unitarizable representation is an irreducible parabolic induction.
Inverting its Satake parameters gives those of the lift: the parameters of $U$ twisted, together with a
segment.

% before: gal-larger-split-satake-minguez
The proof above identifies only the contragredient of the local theta lift, and reaches the lift through
its Satake parameters. Two standard facts identify the lift itself with an induced representation. For
dual pairs of type II the lift has also been computed explicitly by Mínguez, and together with Bernstein's
irreducibility theorem this gives the same induced representation.

% before: gal-larger-lift-enlargement
The weights of the lift at the infinite places and its Satake parameters at the split places are now
known. With attuned characters both take the form required of an enlargement.

% before: gal-larger-pairs
The rank $n+t$ is chosen odd and large enough for the theta lift. The weight bound of the theta lift is
met by giving $\chi_W$ a small infinity type, and the two choices of $\chi_X$ give the two enlargements.

% before: gal-constituent-data
The extraction compares representations of dimensions $n$ and $n+t$, whose Satake normalizations
differ by powers of $q_w$. Stated through the Frobenius eigenvalues and Hodge--Tate multisets themselves,
the four conditions add over direct sums and shift under twists.

% before: gal-constituent-rules
The extraction removes characters from a direct sum and twists by a character. The local conditions
pass through both operations and are invariant under isomorphism.

% before: gal-constituent-characters
The added characters and the twist $\kappa$ act on the Galois side through their $\ell$-adic characters.
These satisfy the local conditions wherever the Hecke character is unramified, and distinct Hecke
characters give distinct Galois characters, so the added characters of the two lifts stay apart.

% before: gal-constituent-peel
The representations attached to the two enlargements have the same characteristic polynomials once each
is multiplied by the added characters of the other. Since the added characters differ, those of the
first enlargement split off from its representation one at a time.

% before: gal-constituent
The representation attached to $U$ is the remaining part, twisted by the character of $\kappa$. Its
construction depends on a place $w_0$ at which the characters are unramified; the representations built
at different places are isomorphic, which gives one representation with the four conditions at every
$d$-split place.

% before: gal-core
The last step passes from definite unitary groups of odd rank to all cuspidal representations in the
self-dual scope. Theta lifts reach even rank, descent from $\operatorname{GL}_a$ to definite unitary
groups reaches cuspidal representations, and base change with gluing reaches all places.
